\documentclass[11pt]{article}
\usepackage[T1]{fontenc}
\usepackage[utf8]{inputenc}
\usepackage{lmodern,amsmath,amssymb,amsthm,mathtools,microtype,booktabs}
\usepackage[margin=27mm]{geometry}
\usepackage[round,authoryear]{natbib}
\usepackage{xurl}
\usepackage[hidelinks]{hyperref}
\newcommand{\workbenchurl}{https://github.com/KokunoYumeto/arithmetic-geometry-complexity-workbench}
\newcommand{\recordurl}{https://github.com/KokunoYumeto/mathematics-workbenches}
\newcommand{\F}{\mathbb F}
\newcommand{\Z}{\mathbb Z}
\newcommand{\Q}{\mathbb Q}
\newcommand{\PP}{\mathbb P}
\newtheorem{theorem}{Theorem}[section]
\newtheorem{proposition}[theorem]{Proposition}
\newtheorem{lemma}[theorem]{Lemma}
\newtheorem{corollary}[theorem]{Corollary}
\theoremstyle{definition}
\newtheorem{definition}[theorem]{Definition}
\title{\#SAT modulo $p$ and Artin--Mazur height,\protect\\after Katz and Stienstra}
\author{}
\date{25 November 2026}
\begin{document}
\maketitle
\begin{abstract}
A companion counting construction, following Katz's coefficient formula,
realizes the number of satisfying assignments modulo a prime as the
Hasse--Witt scalar of a smooth Calabi--Yau hypersurface.
Starting from this construction and Stienstra's logarithm, we identify
its second height test with a count of balanced local assignments,
or equivalently with a count of cyclic orbits.
The resulting height-two decision is modular-counting complete at every
fixed prime, even on the locus where the first scalar vanishes.
At a fixed prime, exact-height decision lies in modular counting with
the requested height supplied in unary and allowed to grow.
Logically equivalent formulas with the same integer satisfying count
can have different height-two status. A binary example also gives
parity-counting completeness at height three. When the prime is part
of the input, even the promised height-two oracle recovers exact
counting through polynomially many queries.
We give exact point-count tests, including the extra precision needed
at two, a finite stopping rule, and polynomial-space computation of
unrestricted exact height. We analyze the existing quasi-Frobenius-splitting
matrix criterion of Kawakami, Takamatsu and Yoshikawa and its implementation
by Batubara, Garzella and Pan; a separate dense implementation takes time
exponential in the number of variables times the logarithm of the prime. A matching complexity
classification for the unrestricted function remains open.
The mathematics and text were produced by GPT-6 Astra (ChatGPT Pro mode);
separate readings by the same model and exact computations are recorded,
and no human has checked the paper.
\end{abstract}

\section{Counting formulas and the height question}
The first coherent Frobenius scalar of a Calabi--Yau hypersurface detects
whether its one-dimensional Artin--Mazur formal group has height one.
It does not determine what happens after that scalar vanishes.
Our question is particularly concrete for the hypersurfaces associated
to Boolean formulas by the companion counting construction
\citep[Theorems~1.1--1.2, Sections~2--3]{Counting}.
That construction gives a smooth geometrically integral hypersurface
$X_\Phi$ over $\mathbb F_p$, with trivial canonical bundle and top
coherent Frobenius scalar $S\bmod p$, where
$S=\#\operatorname{SAT}(\Phi)$ is the number of satisfying truth assignments.
The present paper studies the
additional information in its height.
For a one-dimensional formal group in characteristic $p$, its height
is the integer $h\ge1$ for which the first nonzero term of its
multiplication-by-$p$ series is $aT^{p^h}$, with $a\ne0$; the height
is $\infty$ if that series is zero. The notation $h(X)$ refers to
the Artin--Mazur group associated with the top coherent cohomology
of $X$. Thus a height test distinguishes successive possible first
terms of this intrinsic formal-group operation.

A width-three formula is a conjunction of relations involving at most
three distinct variables. Repeated clauses are separate relations;
unused declared variables are given unary tautologies.
For a positive integer $M$, define $T_M(\Phi)$ as follows.
At each clause choose an ordered list of $M$ satisfying local
assignments, and at each variable choose an ordered list of $M$ bits.
For each incidence of a variable in a clause, require that the total
number of ones in these two lists is $M$.
The number of all such choices is $T_M(\Phi)$.
In particular $T_1(\Phi)=S$. This condition imposes agreement of
one-coordinate marginal counts; it need not impose agreement in
each position. Section~\ref{sec:construction} gives the integral
incidence polynomial, the explicit hypersurface equation, and the
choice of its degree. These definitions specify $X_\Phi$ completely.

Here $\#\mathrm P$ denotes functions counting the accepting branches of
polynomial-time nondeterministic machines. The class
$\mathrm{Mod}_p\mathrm P$ consists of languages accepted when such a
count is nonzero modulo $p$; its binary instance is the parity class
$\oplus\mathrm P$.

\begin{theorem}[The second height stratum]\label{thm:intro-second}
Fix a prime $p$ and use the counting hypersurface $X_\Phi$ defined in
Section~\ref{sec:construction}. If $p\mid S$, then $p\mid T_{p+1}(\Phi)$
and
\[
 h(X_\Phi)=2\quad\Longleftrightarrow\quad
 T_{p+1}(\Phi)/p\not\equiv0\pmod p.
\]
More precisely, simultaneous rotation of the first $p$ positions in
all the lists counted by $T_{p+1}$, leaving the last position fixed,
has $S^2$ fixed points. If $O_p(\Phi)$ is the number of its nontrivial
orbits, then
\[
 T_{p+1}(\Phi)=S^2+pO_p(\Phi),\qquad
 h(X_\Phi)=2\Longleftrightarrow O_p(\Phi)\not\equiv0\pmod p.
\]
For fixed $p$, $O_p$ is a $\#\mathrm P$ function.
Deciding height two on the promise $p\mid S$ is
$\mathrm{Mod}_p\mathrm P$-complete under polynomial-time many-one reductions.
For every prime, two logically equivalent formulas with equal exact
satisfying count, and with the same hypersurface degree, can have
height two and height at least three, respectively.
\end{theorem}

Thus the second arithmetic invariant is a count attached to the
formula's incidence presentation. It cannot be recovered from
$S\bmod p^2$, or even from the integer $S$ alone.
The alternative to height two includes infinite height.
These are completeness statements for counting problems, not
unconditional separations of complexity classes.

\begin{theorem}[A point-count criterion]\label{thm:intro-points}
Let $X\subset\mathbb P^{d-1}_{\mathbb F_p}$ be smooth and geometrically
integral of degree $d\ge4$, and assume its top Artin--Mazur height
is at least two. Put $n=d-2$, $N_r=\#X(\mathbb F_{p^r})$, and
$\tau_r=(-1)^n(N_r-\sum_{i=0}^n p^{ir})$.
Then
\[
 h(X)=2\quad\Longleftrightarrow\quad
 v_p\!\left(\frac{\tau_1^2-\tau_2}{2}\right)=1.
\]
For odd $p$ this is equivalent to $\tau_2/p\not\equiv0\pmod p$,
so $N_2\bmod p^2$ suffices. For $p=2$ it is equivalent to
$(\tau_1^2-\tau_2)/4\not\equiv0\pmod2$, determined by
$N_1\bmod4$ and $N_2\bmod8$.
There exist smooth quintic threefolds with identical $N_1$ and
identical $N_2\bmod4$ but different height-two status.
\end{theorem}

\begin{theorem}[Exact height and its complexity]\label{thm:intro-complexity}
Let $F$ be a division-free arithmetic circuit over $\mathbb F_p$,
with its $d\ge3$ variables explicitly listed, promised to represent
a homogeneous polynomial of degree $d$ defining a smooth hypersurface
$X\subset\mathbb P^{d-1}$. For a homogeneous integral lift put
\[
 a_{p^j}=[(x_1\cdots x_d)^{p^j-1}]\widetilde F^{p^j-1},\qquad
 B_d=\begin{cases}2,&d=3,\\36,&d=4,\\
 \binom{2d-1}{d-1},&d\ge5.\end{cases}
\]
The first index $j\le B_d$ with $a_{p^j}\not\equiv0\pmod{p^j}$
is $h(X)$. If there is no such index, then $h(X)=\infty$.
For every fixed prime $p$ and fixed positive integer $r$, deciding
$h(X)=r$ belongs to $\mathrm{Mod}_p\mathrm P$.
The same membership holds with $r$ supplied in unary and allowed to
grow with the input; the resulting variable-$r$ problem is complete.
The cases $r=1,2$ are complete for that class, and at $p=2$ the
case $r=3$ is $\oplus\mathrm P$-complete even under the promise
$h(X)\ge3$.
The exact-height function, including the output $\infty$, is
computable deterministically with $\operatorname{poly}(L)$ work bits,
where $L$ is the circuit input length, and therefore in
$2^{\operatorname{poly}(L)}$ time. This upper bound also permits a
binary-encoded input prime. A separate dense implementation uses
$2^{O(d\log p)}\operatorname{poly}(s,d,\log p)$ bit operations,
where $s$ is circuit size. Under this uniform
convention the function is $\#\mathrm P$-hard by polynomially many
nonadaptive oracle calls, even on inputs of height at least two. With
a unary-encoded prime, promised height-two decision is oracle-equivalent
to exact counting as stated in Theorem~\ref{thm:uniform-counting}.
The same oracle equivalence holds for arbitrary unary requested heights,
as in Theorem~\ref{thm:unary-height}.
\end{theorem}

The finite coefficient rule determines height in terms of the equation
constructed from $\Phi$. The short replica formula applies only before
the exceptional selections in the stabilizing equation become possible;
we do not extrapolate it to all coefficient indices.
The exact-counting lower bound varies the prime and uses shifted
counting instances and Chinese remaindering. It makes no claim of
$\#\mathrm P$-hardness at one fixed prime.
The polynomial-space upper bound recomputes entries of a fixed-degree
quasi-Frobenius-splitting matrix and treats infinite height explicitly. It does not establish a matching hardness
class for the unbounded exact-height function, while the modular-counting membership for a growing unary height
parameter does not extend to binary height parameters.

\subsection{Classical inputs and the additional arithmetic questions}
Artin and Mazur introduced the formal groups and related their rational
Dieudonn\'e modules to crystalline slopes
\citep[II, Corollaries~4.2--4.3, pp.~109--110; III, Corollaries~3.3--3.4,
pp.~118--119]{AM}. Stienstra gives their logarithm coefficients for
complete intersections \citep[Theorem~1, pp.~907--908]{Stienstra}.
Vlasenko's general congruence theorem already characterizes height
by valuations of those coefficients
\citep[Theorem~2(i), arXiv version~5, p.~3]{Vlasenko2019}.
Our coefficient-to-height lemma supplies the particular leading
$p$-series quotient needed for the incidence count.
The finite stopping bound uses van der Geer and Katsura's
Corollary~2.3, p.~100 \citep{vdGK2003}; the deterministic complexity
bound analyzes the criterion of Kawakami, Takamatsu and Yoshikawa
\citep[Theorem~C, version~3, p.~4; already in version~2, p.~5]{KTY2025}
and the fixed matrix algorithm of Batubara, Garzella and Pan
\citep[Section~3.4 and Algorithm~4, p.~7]{BGP2025}.
Yobuko's comparison identifies its quasi-Frobenius-split height with
Artin--Mazur height \citep[Theorem~4.5, author version pp.~9--10]{Yobuko2019}.
The operator and its iteration are these authors' results. We supply
the uniform resource analysis for circuit inputs, including a finite
matrix cutoff, and the counting interpretation of its coefficients
that yields the unary-height classification. An alternative proof in Appendix~\ref{subsec:height-polynomial-space}
uses Lauder and Wan's finite traces and Hesse, Allender and Barrington's
uniform integer division
\citep[Theorem~28, pp.~600--601]{LauderWan2008}
\citep[Corollary~6.3, author manuscript p.~16]{HAB2002}.
The representation of the coefficients as counts, the examples,
and the promised reductions are proved below.
For computational context, Batubara, Garzella, Huang, and Mellberg
develop practical controlled-reduction implementations and exhibit
explicit Newton strata, including quartic K3 surfaces of specified heights
\citep[Section~1.5, pp.~4--5; Theorem~1.4, p.~5]{BGHM2026}.

The later sections address separate arithmetic questions suggested
by the same complexity programme. They identify the Fermat character
and the correct first variation of coherent Frobenius; state precisely
which succinct motive presentations admit bounded-rank hardness;
prove almost-everywhere stability of orbit-closure witness degrees;
refute the unrestricted class-group all-minors statement;
compute factorial and repeated-root local valuation identities;
and classify affine families of odd theta characteristics.
Their scope is stated with each result. In particular, none gives a
proof-complexity lower bound without an additional covering argument.

\section{Proof overview}
The initial scalar comes from the reciprocal-product generator in the
standard affine cover of projective space. Stienstra's theorem extends
that calculation to an integral logarithm. Comparing its value on the
multiplication-by-$p$ series with $p$ times the logarithm isolates the
first possible nonzero height coefficient. The denominator matters:
after the first scalar vanishes, the next test is a coefficient divided
by $p$, taken modulo $p$.

For the counting equation, balancing the exponents of its two new
coordinates excludes all exceptional selections at the first few
indices. The tag variables then force a prescribed number of copies
of each incidence factor. This gives the balanced-list count $T_M$.
A cyclic action on the lists explains its divisibility, and constant
formulas with a nonzero second quotient turn the identity into a
promise-preserving hardness reduction. The binary third-height example
uses an exact finite binomial sum.

The point-count test follows from a different description of the same
formal group. At height two there are precisely two slopes equal to
one half. Their product is the unique term of valuation one in the
second elementary symmetric coefficient. Higher and infinite heights
force larger valuation. Newton's identity then yields the stated
point-count precisions. Finally, a finite cohomological height bound
makes the logarithm rule terminate. For polynomial space, the quasi-Frobenius-splitting criterion
keeps every iterate in the same homogeneous polynomial space. Its
matrix entries use only scalar residues modulo the prime and its square.
Cayley--Hamilton supplies a finite stopping rule. Recursive coefficient
convolutions and matrix powering compute any requested entry with a
polynomial stack, while a separate dense implementation gives the
stated sharper time bound. Finally, varying
the prime, shifting satisfying counts and applying Chinese remaindering
recovers exact counting from the promised height-two oracle.

\section{The first scalar and the point-count congruence}
This section fixes the normalization of the scalar before asking for higher-height information.
For a degree-$d$ equation $F$ in $d$ variables, define its Hasse coefficient by
\[
 h_p(F)=[(x_1\cdots x_d)^{p-1}]F^{p-1}\in\mathbb F_p.
\]
Here $[m]A$ denotes the coefficient of the monomial $m$ in a polynomial $A$.

\begin{proposition}\label{prop:scalar}
For a degree-$d$ hypersurface as above, $H^{d-2}(X,\mathcal O_X)$ is one-dimensional over $\F_p$. On the basis induced by the reciprocal monomial $(x_1\cdots x_d)^{-1}$, absolute Frobenius acts by $h_p(F)$. Moreover
\[
 \#X(\F_p)\equiv1+(-1)^d h_p(F)\pmod p.
\]
\end{proposition}
\begin{proof}
Multiplication by $F$ gives the exact sequence
\[
 0\longrightarrow\mathcal O_{\PP^{d-1}}(-d)
 \xrightarrow{\ F\ }\mathcal O_{\PP^{d-1}}
 \longrightarrow\mathcal O_X\longrightarrow0.
\]
On the standard cover $x_i\ne0$, the top cohomology of $\mathcal O(-d)$ is the degree-$-d$ part of the Laurent polynomial ring modulo the sum of subrings in which at least one exponent is nonnegative. Its only surviving monomial is $(x_1\cdots x_d)^{-1}$, so that cohomology is a line. The same cover gives zero intermediate cohomology of the two projective-space sheaves in the positions needed here. For $d=3$ this is the usual vanishing of $H^1(\PP^2,\mathcal O)$ and $H^2(\PP^2,\mathcal O)$; the same monomial computation proves it for every $d$. The connecting homomorphism therefore identifies $H^{d-2}(X,\mathcal O_X)$ with this line.

The map on the left sheaf which makes Frobenius commute with multiplication by $F$ is $a\mapsto F^{p-1}a^p$, since $F(F^{p-1}a^p)=(Fa)^p$. Applied to the indicated generator it gives
$F^{p-1}/(x_1\cdots x_d)^p$. A monomial of $F^{p-1}$ survives only if every exponent is at most $p-1$. Their sum is $d(p-1)$, so all must equal $p-1$. The surviving coefficient is precisely $h_p(F)$.

For the point count, let $A$ be the number of affine zeros of $F$. Fermat's identity on $\F_p$ gives
$A\equiv-\sum_{x\in\F_p^d}F(x)^{p-1}\pmod p$.
For an exponent $m$, $\sum_{a\in\F_p}a^m$ is zero when $m=0$, zero when $m>0$ is not divisible by $p-1$, and $-1$ otherwise. Since the total degree is $d(p-1)$, exactly the coefficient defining $h_p(F)$ survives the product of these sums. Thus $A\equiv(-1)^{d+1}h_p(F)$.
The affine cone satisfies $A=1+(p-1)\#X(\F_p)$, hence the claimed congruence. It has a plus sign when $d$ is even, and also in characteristic two. Keeping track of this parity is necessary when a counting construction identifies its scalar with a prescribed residue.
\end{proof}

\section{The counting equation and its integral coefficients}\label{sec:construction}
This section recalls the counting construction with the normalization needed for higher coefficients. The construction is due to the companion counting paper \citep[Theorems~1.1--1.2, Sections~2--3]{Counting}; the calculations below check the integral information that its first Hasse coefficient does not record.

Write a width-three formula in conjunctive normal form (CNF) as a conjunction of local relations $f$, with variable scopes $I_f$ that are sets of distinct variables. Different clauses may have the same scope. A declared variable absent from the formula is given one unary tautological relation. If the incidence set is empty, add a fresh variable with a unary clause forcing it to zero; this makes the incidence set nonempty without changing the count. Write $\mathcal I=\{(f,i):i\in I_f\}$ for that set. For each incidence $e=(f,i)$ introduce $w_{e,0},w_{e,1}$. Put
\[
 P_f=\sum_{a\models f}\prod_{i\in I_f}w_{(f,i),a_i},\qquad
 Q_i=\prod_{f:i\in I_f}w_{(f,i),0}+
     \prod_{f:i\in I_f}w_{(f,i),1},
\]
\[
 P=\prod_fP_f\prod_iQ_i,\qquad
 Z=\prod_e w_{e,0}w_{e,1},\qquad b=2|\mathcal I|.
\]
The coefficient $[Z]P$ is the number $S$ of satisfying assignments: the variable factor chooses a consistent label, and every neighboring clause factor must choose its complement. This gives a bijection, including the multiplicities of equal monomials. In particular $P$ is an integral, homogeneous, positive circuit of degree $b$, and every individual variable has degree at most two.

Introduce $p-1$ tags and set $u=b+p-1$ and
\[
 G=-t_1^{p-1}P+\sum_{j=2}^{p-1}t_j^{p-1}Z.
\]
The empty sum is zero. This is degree $u$ in $u$ variables. Choose the least power $Q=p^r$ with $Q-1>u$ and $Q-1\geq p$, put $m=Q-1$ and $d=Q+1$, and relabel the variables of $G$ as $x_1,\ldots,x_u$. Define integral polynomials
\[
 A=G\prod_{i=u+1}^{m}x_i+x_1^m,\qquad
 F=yzA+y^{Q+1}+\sum_{i=1}^{m}x_i^{Q+1}-x_1z^Q.
\]
Reduction modulo $p$ defines the counting hypersurface $X_\Phi=V(F)\subseteq\PP^{d-1}$.

For clarity, its smoothness is uniform in every coefficient of $G$. In characteristic $p$,
\[
 F_y=zA+y^Q,\quad F_z=yA,\quad yF_y-zF_z=y^{Q+1}.
\]
A common zero of the partial derivatives has $y=0$, then $x_i=0$ for $i>1$ and $x_1^Q=z^Q$. At a projective point both $x_1,z$ are nonzero. The padding makes $A(x_1,0,\ldots,0)=x_1^m$, so $F_y=zx_1^m\ne0$, a contradiction. Adjunction and the hypersurface cohomology sequence give trivial canonical sheaf, a one-dimensional top coherent group, and zero intermediate coherent groups. After every extension of the base field, the same sequence gives $H^0(\mathcal O)$ equal to that field; the smooth geometric hypersurface is therefore connected and integral. These arguments also establish the hypotheses used for the formal group.

Let
\[
 a_N(F)=[(x_1\cdots x_myz)^{N-1}]F^{N-1},\qquad
 T_M(\Phi)=[Z^M]P^M.
\]
\begin{proposition}\label{prop:early-counting-coefficients}
If $p^k-1<Q(Q+1)$ and $M=(p^k-1)/(p-1)$, then over the integers
\[
 a_{p^k}(F)=(-1)^M
 \frac{(p^k-1)!}{(M!)^{p-1}}\,T_M(\Phi).
\]
The displayed multiplier is a $p$-adic unit congruent to one modulo $p$.
\end{proposition}
\begin{proof}
Put $n=p^k-1$. In a contributing term let $a,c,t$ count the selections of $y^{Q+1}$, $-x_1z^Q$, and $yzA$, respectively. The $y,z$ exponents give
$(Q+1)a+t=n=Qc+t$, so $a=Qj,c=(Q+1)j$ for an integer $j\geq0$, and $t=n-Q(Q+1)j$. The assumed inequality forces $j=0$. All $n$ factors therefore select $yzA$. Any padding variable has target exponent $n$, and occurs only in the first summand of $A$, with exponent one. Consequently every factor selects that summand. The coefficient reduces exactly to the diagonal coefficient of $G^n$.
Each tag then forces precisely $M$ selections of its summand. Removing the $p-2$ copies of the $Z$ summand leaves $[Z^M]P^M$, with the asserted sign and multinomial coefficient. The base-$p$ digits of $M$ are $k$ ones. Adding $p-1$ copies produces $k$ digits equal to $p-1$ without a carry, so the multinomial is a unit. Its residue is $((p-1)!)^k=(-1)^k$ by the elementary factorial digit congruence. For odd $p$, $M\equiv k\pmod2$, so the sign cancels this residue. At $p=2$ signs agree modulo $p$.
\end{proof}


\section{Stienstra coefficients and Artin--Mazur height}
\label{sec:classical-height}
Let $p$ be any prime, including $2$, and let $F$ be a homogeneous
polynomial of degree $d$ in $d\ge3$ variables over $\mathbb F_p$,
defining a smooth projective hypersurface $X$.
Choose a homogeneous lift $\widetilde F\in\mathbb Z_p[x_1,\ldots,x_d]$
and put $z=x_1\cdots x_d$. Define
\begin{equation}
 a_m=a_m(\widetilde F)=[z^{m-1}]\widetilde F^{\,m-1}\quad(m\ge1).
 \label{eq:H1}
\end{equation}
Thus $a_1=1$ and $a_p\bmod p=h_p(F)$.

\subsection{The source theorem and its hypotheses}
Stienstra's Theorem~1, printed p.~907, with its base-change statement
on p.~908 \citep{Stienstra}, gives an integral one-dimensional formal
group law for the Artin--Mazur group of $X$ with logarithm
\begin{equation}
 \ell(T)=\sum_{m\ge1}\frac{a_m}{m}T^m.                 \label{eq:H2}
\end{equation}
His general theorem concerns a homogeneous regular sequence
$F_1,\ldots,F_r$ over a noetherian ring $K$, with flat projective
intersection, and degrees $d_i\ge \sum_i d_i-N\ge1$ in $\mathbb P^N$.
The logarithm statement additionally assumes that $K$ is flat over
$\mathbb Z$. In our specialization $r=1$, $N=d-1$, and the index set in
his theorem is the singleton $(1,\ldots,1)$, which gives exactly \eqref{eq:H1}.
A primitive hypersurface equation over $\mathbb Z_p$ gives a flat
hypersurface, so these hypotheses hold for any lift of the nonzero
equation $F$.

Artin--Mazur prove representability for projective complete intersections
of dimension at least two in the special case following
II, Corollary~(4.2), pp.~109--110; II, Corollary~(4.3), p.~110 identifies
the Dieudonn\'e module with Witt-vector cohomology
\citep{AM}. No crystalline comparison is needed for the
height-one or height-two computation below.

\subsection{An elementary logarithm lemma}
Vlasenko's general congruence theorem already determines height from
these valuations: in her notation $b_{p^j-1}=a_{p^j}$, finite height
$h$ gives $v_p(a_{p^j})\ge j-\lfloor j/h\rfloor$, with equality when
$h\mid j$, and infinite height gives $v_p(a_{p^j})\ge j$
\citep[Theorem~2(i), arXiv version~5, p.~3]{Vlasenko2019}.
The following proof also identifies the leading multiplication-by-$p$
coefficient in the chosen coordinate, which fixes the quotient used
in the counting interpretation.
\begin{lemma}
\label{lem:logarithm-height}
Let $\mathcal F$ be a one-dimensional commutative formal group law over
$\mathbb Z_p$ whose logarithm has the form
$\ell(T)=\sum_{m\ge1}a_mT^m/m$, with $a_m\in\mathbb Z_p$ and $a_1=1$.
Let $h$ be the height of its reduction, and write
$B(T)=[p]_{\mathcal F}(T)=\sum_{j\ge1}b_jT^j$.
For every $r\ge1$ for which $h\ge r$,
\begin{equation}
 a_{p^r}\in p^{r-1}\mathbb Z_p,\qquad
 b_{p^r}\equiv\frac{a_{p^r}}{p^{r-1}}\pmod p.         \label{eq:H3}
\end{equation}
Consequently, under that hypothesis,
\begin{equation}
 h=r\quad\Longleftrightarrow\quad
 \frac{a_{p^r}}{p^{r-1}}\not\equiv0\pmod p.           \label{eq:H4}
\end{equation}
\end{lemma}
\begin{proof}
We first recall the elementary power-of-$p$ property, including its
reason. If $g$ is a homomorphism of one-dimensional formal group laws
$\mathcal A,\mathcal B$ over a field of characteristic $p$, differentiation
of the homomorphism identity in its second variable at zero gives
\[
 g'(T)\,\partial_2\mathcal A(T,0)
 =\partial_2\mathcal B(g(T),0)\,g'(0).
\]
Both displayed partial derivatives are units. Thus $g'(0)=0$ implies
$g'(T)=0$ and $g(T)=g_1(T^p)$. The original identity and the injectivity
of substitution $(U,V)\mapsto(U^p,V^p)$ show that $g_1$ is a
homomorphism from the Frobenius twist of $\mathcal A$ to $\mathcal B$.
Iteration proves that the first nonzero term of a nonzero $p$-series
has degree $p^h$. It also proves that if its coefficient at $p^r$
vanishes after $r$ such factorizations, one can factor once more.
The zero $p$-series has height infinity.

Suppose $h\ge r$, and set $M=p^r$. Then $b_j\in p\mathbb Z_p$ for
$1\le j<M$. Compare the coefficient of $T^M$ in
$\ell(B(T))=p\ell(T)$:
\begin{equation}
 b_M+\sum_{j=2}^{M}\frac{a_j}{j}[T^M]B(T)^j
       =\frac{a_M}{p^{r-1}}.                        \label{eq:H5}
\end{equation}
Every term contributing to $[T^M]B^j$ uses $j$ coefficients $b_i$
with $i\le M-j+1<M$, and is therefore divisible by $p^j$.
Since $j-v_p(j)\ge1$ for $j\ge2$, each summand in \eqref{eq:H5} is an
integral multiple of $p$. This proves both the integrality assertion
and \eqref{eq:H3}. The power-of-$p$ property then gives \eqref{eq:H4}.
There is no exception at $p=2$.
\end{proof}

Applied to \eqref{eq:H2}, this proves the two tests
\begin{equation}
 h(X)=1\iff a_p\not\equiv0\pmod p,                  \label{eq:H6}
\end{equation}
and, on the locus $a_p\equiv0\pmod p$,
\begin{equation}
 h(X)=2\iff a_{p^2}/p\not\equiv0\pmod p.            \label{eq:H7}
\end{equation}
Here $a_{p^2}$ means the coefficient in $\widetilde F^{p^2-1}$,
not the square of $a_p$. Formula \eqref{eq:H7} requires that coefficient
modulo $p^2$.

\subsection{Lift independence and a universal second equation}
For a degree-$d$ polynomial $A$ in $d$ variables, write
$A_1(A)=[z^{p-1}]A^{p-1}$ and
$A_2(A)=[z^{p^2-1}]A^{p^2-1}$.
In characteristic $p$,
\begin{equation}
 A_2(A)=A_1(A)^{p+1}.                              \label{eq:H8}
\end{equation}
Indeed, factor $A^{p^2-1}=A^{p-1}(A^{p-1})^p$.
A contributing exponent in the first factor is congruent to $p-1$
in every coordinate. Each coordinate is at least $p-1$, while the
total degree is $d(p-1)$, forcing every coordinate to equal $p-1$.

Consequently, if the coefficients of $A$ are independent integer
variables, the polynomial
\begin{equation}
 H_2(A)=\frac{A_2(A)-A_1(A)^{p+1}}p                \label{eq:H9}
\end{equation}
has integer coefficients. On $H_1=A_1\bmod p=0$, its reduction agrees
with $a_{p^2}/p\bmod p$, since $a_p^{p+1}/p$ is divisible by $p$.
Thus height at least three is cut out, on the smooth parameter locus,
by $H_1=H_2=0$, and height exactly two by $H_1=0$, $H_2\ne0$.
This states equations for the loci; it makes no assertion about their
codimension, reducedness, smoothness, or nonemptiness.

There is also a direct check of the quotient's independence of the
chosen lift. Write $\widetilde F'=\widetilde F+pH$. Modulo $p^2$,
\[
 a_{p^2}(\widetilde F')-a_{p^2}(\widetilde F)
 \equiv p(p^2-1)[z^{p^2-1}]H\widetilde F^{p^2-2}.
\]
The same exponent argument, now applied to
$HF^{p-2}(F^{p-1})^p$, gives
\begin{equation}
 \frac{a_{p^2}(\widetilde F')-a_{p^2}(\widetilde F)}p
 \equiv
 -h_p(F)^p\,[z^{p-1}]HF^{p-2}\pmod p.             \label{eq:H10}
\end{equation}
It vanishes when $h_p(F)=0$. This proof includes $p=2$, where
$F^{p-2}=1$. Hence \eqref{eq:H7} is an intrinsic test of the special fibre,
not a choice of lift.

\subsection{A precise $p$-typical recurrence}
For clarity, one can obtain a $p$-typical logarithm without assuming
that the original logarithm \eqref{eq:H2} already is $p$-typical.
For a prime $q\ne p$ and an integral curve $\gamma(T)$ of $\mathcal F$,
apply the operator
\begin{equation}
 \gamma(T)\longmapsto
 \gamma(T)-_{\mathcal F}[1/q]_{\mathcal F}
       \Bigl(\mathop{\sum\nolimits^{\mathcal F}}_{\zeta^q=1}
                         \gamma(\zeta T)\Bigr).    \label{eq:H11}
\end{equation}
The sum can be formed over a finite unramified extension containing
the $q$th roots of unity; it is invariant and descends.
The inverse of $[q]_{\mathcal F}$ is integral because its tangent
coefficient $q$ is a unit. On logarithms, \eqref{eq:H11} removes precisely
the degrees divisible by $q$. Starting with $\gamma(T)=T$ and taking
the product over all primes $q\ne p$ is well defined degree by degree,
and gives an integral strict coordinate change $c(T)=T+O(T^2)$.
The logarithm in that coordinate is
\begin{equation}
 L(T)=\sum_{r\ge0}\lambda_rT^{p^r},\qquad
 \lambda_0=1,\quad \lambda_r=a_{p^r}/p^r.          \label{eq:H12}
\end{equation}

Let $\mathcal F_p$ be this formal group law. There are unique integral
parameters $w_r$, with $w_0=p$, such that
\begin{equation}
 [p]_{\mathcal F_p}(T)
 =\mathop{\sum\nolimits^{\mathcal F_p}}_{r\ge0}w_rT^{p^r}.
                                                               \label{eq:H13}
\end{equation}
To construct them, successively subtract the leading term of the
remaining curve using formal-group subtraction. Its logarithm has only
$p$-power degrees, so after the degree $p^r$ term is removed, the next
possible leading degree is $p^{r+1}$. The removed leading coefficient
is integral, which proves the assertion.

Taking logarithms of \eqref{eq:H13} gives the exact recurrence
\begin{equation}
 p\lambda_n=\sum_{i=0}^n\lambda_{n-i}w_i^{p^{n-i}},
 \qquad
 w_n=(p-p^{p^n})\lambda_n
        -\sum_{i=1}^{n-1}\lambda_{n-i}w_i^{p^{n-i}}.              \label{eq:H14}
\end{equation}
In particular,
\begin{equation}
 w_1=(1-p^{p-1})a_p,\qquad
 w_2=(1-p^{p^2-1})\frac{a_{p^2}}p
                   -\frac{a_p}{p}w_1^p.                       \label{eq:H15}
\end{equation}
The sum in \eqref{eq:H13} is formal-group addition, not ordinary addition.
Dropping the $p^{p^n}\lambda_n$ term while retaining the same
parameters would give an incorrect exact recurrence.
Modulo $p$, $w_1=H_1$ and $w_2=H_2$; the latter follows from
$(1-p^{p-1})^p\equiv1\pmod{p^2}$.
On $H_1=0$, \eqref{eq:H15} reduces to \eqref{eq:H7}.
The direct proof \eqref{eq:H5} establishes the height tests without requiring
this choice of coordinate.

\subsection{Finite tests and the crystalline hypothesis boundary}
For a prescribed cutoff $R$, formulas \eqref{eq:H3}--\eqref{eq:H4} give a finite procedure:
test $a_p\bmod p$, then, on its zero locus, test $a_{p^2}\bmod p^2$
after division by $p$, and continue through
$a_{p^R}/p^{R-1}\bmod p$.
The result is the exact height if it is at most $R$, or the statement
$h>R$. A finite cutoff alone does not identify infinite height.

For the smooth projective complete intersections here, one can choose
a smooth proper lift to $\mathbb Z_p$. Indeed, a homogeneous lift of
the nonzero equation gives a flat projective hypersurface. Its special
fibre is smooth, and properness makes the image of the closed
nonsmooth locus closed in $\operatorname{Spec}\mathbb Z_p$.
That image misses the closed point and must therefore be empty.
Artin--Mazur III,
Corollary~(1.9)(a), pp.~114--115, and III, Corollary~3.3 with
Remark~2, pp.~118--119, verify the hypotheses needed to identify the
Dieudonn\'e module, after inverting $p$, with the slope-$<1$ part of middle crystalline
cohomology \citep{AM}. In finite height its rank is the
height, so the middle crystalline rank supplies a finite upper bound.
Invoking that comparison for a general variety without the required
representability and lifting hypotheses would not be justified by
this citation.

\section{Replica counts and the second height stratum}
The incidence factorization converts the logarithm coefficient into an exact finite counting problem.
For a clause $f$ let $\mathcal A_f\subseteq\{0,1\}^{I_f}$ be its
satisfying assignments, and define
\[
 C_{f,M}(r)=
 [\prod_{i\in I_f}z_i^{r_i}]
 \left(\sum_{a\in\mathcal A_f}\prod_{i\in I_f}z_i^{a_i}\right)^M.
\]
This counts ordered lists of $M$ satisfying local assignments with
the specified one-coordinate marginals.  Expanding the variable
factors $Q_i^M$ gives the exact formula
\begin{equation}\label{eq:replica-marginals}
 T_M(\Phi)=\sum_{r\in\{0,\ldots,M\}^{V(\Phi)}}
       \prod_i\binom Mr_i\prod_f C_{f,M}(r|_{I_f}).
\end{equation}
Indeed, choosing the zero label $r_i$ times in $Q_i^M$ forces the
adjacent clause lists to choose the one label $r_i$ times.  This is a
marginal-consistency count, not a count of $M$ globally consistent
satisfying assignments.

\begin{proposition}[The second coefficient as a count of orbits]
Write $S=\#\mathrm{SAT}(\Phi)$.  There is a nonnegative integer
$O_p(\Phi)$, counting explicitly specified orbits, such that
\begin{equation}\label{eq:replica-orbits}
 T_{p+1}(\Phi)=S^2+pO_p(\Phi).
\end{equation}
For fixed $p$, $O_p$ is a \#P function.
\end{proposition}
\begin{proof}
Use the ordered lists counted by \eqref{eq:replica-marginals}, including
the $p+1$ choices at every variable factor.  Simultaneously rotate the
first $p$ positions in every list and leave the last position fixed.
This action preserves incidence balance.  A fixed configuration has
constant first-$p$ choices at each factor.  At an incidence the one-label
balance has the form
\[
 p(a+c)+(b+e)=p+1,\qquad a,b,c,e\in\{0,1\}.
\]
For every prime, including two, this forces $a+c=b+e=1$.
The first block and the final position therefore specify two global
satisfying assignments, and conversely each ordered pair gives one
fixed configuration.  There are $S^2$ fixed points.  Every other orbit
has size $p$, giving \eqref{eq:replica-orbits}.
For fixed $p$, a nondeterministic polynomial-time machine guesses one
configuration, checks all local constraints and balances, and accepts
exactly when it is the lexicographically least member of a nonconstant
rotation orbit.  This counts $O_p$.
\end{proof}



\begin{theorem}[Exact height-two criterion for the counting construction]
For every fixed prime $p$, on the promise $S\equiv0\pmod p$,
\[
 h(X_\Phi)=2
 \quad\Longleftrightarrow\quad
 \frac{T_{p+1}(\Phi)}p\not\equiv0\pmod p
 \quad\Longleftrightarrow\quad O_p(\Phi)\not\equiv0\pmod p.
\]
If these residues vanish, the height is at least three, with infinite
height included.  The promised decision problem is
$\mathrm{Mod}_p\mathrm P$-complete under polynomial-time many-one
reductions.
\end{theorem}
\begin{proof}
The coefficient formula, its unit congruent to one, and the preceding
height test give the first equivalence.  Since $p\mid S$,
$S^2/p\equiv0\pmod p$, and \eqref{eq:replica-orbits} gives the second.
The orbit count proves membership.  The gadgets constructed below
satisfy $p\mid S_\Psi$ and $v_p(T_{p+1}(\Psi))=1$ for every prime.
For an arbitrary input $\Phi$, take its conjunction with $\Psi$ on
disjoint variable sets.  The satisfying count and replica coefficients
both multiply.  Its satisfying count is always divisible by $p$, and
\[
 \frac{T_{p+1}(\Phi\wedge\Psi)}p
 \equiv S_\Phi^2\frac{T_{p+1}(\Psi)}p\pmod p.
\]
This is nonzero exactly when $S_\Phi\not\equiv0\pmod p$.
The gadget is constant size for fixed $p$, and the counting
construction has polynomial circuit size, completing the reduction.
\end{proof}

\subsection{Uniform gadgets and equal-count examples}
For an odd prime let $n=p-1$ and let $\Theta_n$ consist of all binary
clauses $\neg x_i\vee\neg x_j$, $1\leq i<j\leq n$.
It has exactly $n+1=p$ satisfying assignments.  Let $\Theta_{n,j}$
add $j$ copies of the unary tautological relation $x_1\vee\neg x_1$.
Repeated clauses are retained as separate factors, exactly as in the
input construction.

\begin{proposition}[A uniform odd-prime calculation]
For every odd prime,
\begin{equation}\label{eq:generic-height-gadget}
 T_{p+1}(\Theta_{p-1,j})\equiv(-6-2j)p\pmod {p^2}.
\end{equation}
Consequently $\Theta_{p-1,p-3}$ has height at least three and
$\Theta_{p-1,p-2}$ has height two after stabilization, although both
have exactly $p$ satisfying assignments and are logically equivalent.
\end{proposition}
\begin{proof}
Put $M=p+1$ and $B=p+1$.  An intermediate marginal
$2\leq r_i\leq p-1$ makes its variable binomial divisible by $p$.
An incident clause histogram is also divisible by $p$: modulo $p$,
its generating polynomial to the power $p+1$ is the product of that
polynomial at $z^p$ and at $z$, whose possible individual exponents
are only $0,1,p,p+1$.  Hence every term with an intermediate marginal
vanishes modulo $p^2$.

For a pair clause the local coefficient is
\[
 C(r_i,r_j)=
 \begin{cases}
 M!/(r_i!r_j!(M-r_i-r_j)!),&r_i+r_j\leq M,\\
 0,&r_i+r_j>M.
 \end{cases}
\]
At boundary marginals, two entries equal to $1$ supply a factor
$pB$.  Therefore three or more entries equal to $1$ already vanish
modulo $p^2$.  At most one entry can equal $p$; an entry equal to
$p+1$ forces all others to zero.  The complete list of surviving
types for $j=0$ is
\[
\begin{array}{c|c|c}
\text{type}&\text{multiplicity}&\text{weight modulo }p^2\\\hline
\text{all zero}&1&1\\
\text{one }1&n&B^n\\
\text{two }1&\binom n2&p\\
\text{one }p&n&B^n\\
\text{one }p\text{ and one }1&n(n-1)&B^{2n-1}\\
\text{one }p\text{ and two }1&n\binom{n-1}2&p\\
\text{one }p+1&n&1.
\end{array}
\]
Using $B^a\equiv1+ap\pmod{p^2}$, their sum is
\[
 (n+1)^2+p\frac{n(5n^2-4n+3)}2.
\]
For $n=p-1$ its correction is $-6p$ modulo $p^2$.

One additional unary tautology multiplies the term with marginal
$r_1$ by $\binom M{r_1}$.  Its change modulo $p^2$ is zero at
$r_1=0,M$, equals $p$ times the old term at $r_1=1,p$, and
intermediate terms already vanish.  Modulo $p$, boundary configurations
are uniquely ordered pairs of satisfying assignments.  Exactly
$2s_1(S-s_1)=2n$ such pairs disagree at $x_1$, where $s_1=1$.
Thus each added tautology changes the coefficient by $2np=-2p$
modulo $p^2$.  This proves \eqref{eq:generic-height-gadget}.
The two chosen values of $j$ give residues $0$ and $-2p$.
\end{proof}

For $p\geq5$ the two displayed formulas have respectively
$u=2p^2-3p-3$ and $u=2p^2-3p-1$, so the original least-power rule
selects $Q=p^3$ for both.  They thus give equal-count, different-height
examples even in the same ambient dimension.  For $p=3$, adjoining
one forced dummy variable to both leaves $S,T_M$ unchanged and makes
the original choice $Q=27$ for both.

For $p=2$ take
\[
 \Psi_2=(\neg x_1\vee\neg x_2)\wedge
        (\neg x_1\vee\neg x_3)\wedge
        (\neg x_1\vee x_2\vee x_3).
\]
It is logically equivalent to $\neg x_1$, with $x_2,x_3$ free, so
$S=4$.  Formula \eqref{eq:replica-marginals} gives
$T_3(\Psi_2)=47362\equiv2\pmod4$, independently checked by expanding
the incidence polynomial and summing triples of its monomials.
It is therefore the required height-two gadget.  As an equal-count
comparison take $\neg x_1$, four copies of $x_1\vee\neg x_1$, and
one tautology in each of $x_2,x_3$.  Its forced marginal at $x_1$
leaves
\[
 T_3=\binom63^2=400\equiv0\pmod4.
\]
Both formulas have seven incidences, $u=15$, and original $Q=32$.
They define smooth hypersurfaces of degree $33$ in $33$ variables,
with the same exact satisfying count and heights respectively two
and at least three.  Thus even the full integer $S$, and hence
$S\bmod p^2$, does not determine the height for this construction.


\section{An exact third-height stratum in characteristic two}
A constant Boolean formula provides the next nonzero height quotient.
\begin{proposition}[An exact height-three binary gadget]
Let
\[
 \Psi_3=(\neg x_1\vee\neg x_2)\wedge
 (\neg x_1\vee\neg x_3)\wedge
 (\neg x_1\vee\neg x_2\vee\neg x_3)\wedge
 (x_1\vee\neg x_2\vee\neg x_3).
\]
Then $S_{\Psi_3}=4$ and its incidence replica coefficients are
\[
 T_3=1002460\equiv0\pmod4,\qquad
 T_7=1382535756736336500\equiv4\pmod8.
\]
Its counting hypersurface in characteristic two has height exactly
three.  It has ten incidences, $u=21$, and original stabilization
parameter $Q=32$.
\end{proposition}
\begin{proof}
The Boolean function is ``at most one of $x_1,x_2,x_3$ is true'',
so its satisfying count is four.  Here is a finite integer formula
and evaluation certificate for the two further coefficients.
Set $\binom nk=0$ outside $0\leq k\leq n$, and put
\[
 A_M(a,b)=
 \begin{cases}
 \dfrac{M!}{a!b!(M-a-b)!},&a+b\leq M,\\
 0,&a+b>M,
 \end{cases}
\]
and, for $f\in\{0,1\}^3$, put
\[
 C_{M,f}(a,b,c)=\sum_{j=0}^M(-1)^j\binom Mj
 \binom{M-j}{a-jf_1}\binom{M-j}{b-jf_2}
 \binom{M-j}{c-jf_3}.
\]
The pair clause permits exactly the three local assignments
$00,10,01$, so its $M$-replica coefficient is $A_M$.
A ternary clause with unique forbidden assignment $f$ has local
polynomial
$(1+z_1)(1+z_2)(1+z_3)-z_1^{f_1}z_2^{f_2}z_3^{f_3}$.
The binomial expansion of its $M$th power gives $C_{M,f}$.
Multiplying the two pair clauses, the two ternary clauses, and
the variable marginal counts therefore gives the exact formula
\[
 T_M=\sum_{a,b,c=0}^M
 \binom Ma\binom Mb\binom Mc\,
 A_M(a,b)A_M(a,c)C_{M,(1,1,1)}(a,b,c)
 C_{M,(0,1,1)}(a,b,c).
\]
For transparency, grouping this finite sum by $a$ gives the
following integer certificate; the absent $M=3$ entries are zero.
\[
\begin{array}{c|r|r}
 a&M=3&M=7\\\hline
 0&9075&665953307379\\
 1&857547&5184197281860583\\
 2&135837&377589783130120323\\
 3&1&909264005602246875\\
 4&0&90276840331103125\\
 5&0&220262726359053\\
 6&0&1711339161\\
 7&0&1\\\hline
 \text{sum}&1002460&1382535756736336500
\end{array}
\]
Summing the displayed entries gives the asserted residues. The early-stabilization
formula gives $a_2=-S$, $a_4=-T_3$, and $a_8=-T_7$.
The first two height tests vanish and $a_8/4$ is odd, proving the claim.
\end{proof}

\begin{lemma}[The diagonal replica residue]
For the incidence polynomial of any input formula and every prime $p$,
\[
 T_{1+p+\cdots+p^{k-1}}\equiv S^k\pmod p.
\]
\end{lemma}
\begin{proof}
Modulo $p$, the relevant power of $P$ is
$\prod_{j=0}^{k-1}P(z^{p^j})$.  Every individual incidence variable
has exponent at most two in a monomial of $P$.  To obtain target
exponent $1+p+\cdots+p^{k-1}$, reduction modulo $p$ forces the
lowest exponent to be one; this holds also for $p=2$, since the
possibilities are $0,1,2$.  There is no carry.  Induction on the
digits forces exponent one in every factor and every variable.
The target coefficient is therefore the product of $k$ copies of
$[Z]P=S$, modulo $p$.
\end{proof}

\begin{corollary}[The next binary height decision]
On the promise that the counting hypersurface has height at least
three, deciding height three versus height at least four, including
infinity, is $\oplus\mathrm P$-complete.
\end{corollary}
\begin{proof}
For hardness conjoin an arbitrary input $\Phi$ with $\Psi_3$ on
disjoint variables.  Its satisfying count is even and its $T_3$ is
divisible by four, so its height is at least three.  Its height is
three precisely when
\[
 T_7(\Phi)T_7(\Psi_3)/4\equiv T_7(\Phi)
 \equiv S_\Phi\pmod2
\]
is nonzero.  All three early coefficient identities apply to the
combined construction.  For membership $T_7$ counts polynomial-size
configurations, since seven is fixed.  Guess four distinct such
configurations in increasing lexicographic order.  The number of
accepting guesses is $\binom{T_7}{4}$, whose parity is the binary
digit of $T_7$ in position two.  On the promise, this is exactly
$T_7/4\bmod2$, the third height test.  Thus it is a \#P parity.
\end{proof}


\section{Height two from two point counts}
\label{sec:point-count-height-two}
Let $X/\mathbb F_p$ be a smooth geometrically integral hypersurface
of degree $d\ge4$ in $\mathbb P^{d-1}$, and put $n=d-2$.
Its top Artin--Mazur formal group is one-dimensional.
Write $N_r=\#X(\mathbb F_{p^r})$ and
\begin{equation}
 \tau_r=(-1)^n\left(N_r-\sum_{i=0}^{n}p^{ir}\right).
                                                               \label{eq:PC1}
\end{equation}
The Frobenius convention is the one in the point-count trace formula.

\subsection{Classical inputs with their hypotheses}
Weak Lefschetz and duality identify all cohomology outside the middle
with the corresponding projective-space classes. One can recover this
particular weak Lefschetz statement from affine vanishing
\citep[Theorem~15.1, p.~105]{MilneLEC2013} and Poincar\'e duality
\citep[Theorem~24.1, p.~144]{MilneLEC2013}.
Indeed, the complement $U=\mathbb P^{n+1}\setminus X$ is smooth affine
of dimension $n+1$. Duality and affine vanishing give
$H_c^j(U,\mathbb Q_\ell)=0$ for $j<n+1$.
The compact-support exact sequence therefore makes
$H^i(\mathbb P^{n+1})\to H^i(X)$ an isomorphism for $i<n$ and
an injection for $i=n$. Duality on $X$ treats the degrees above $n$.
When $n$ is even, the middle ambient class has nonzero self-pairing
equal to the degree, so its orthogonal complement is Frobenius stable.
When $n$ is odd there is no middle ambient class.

Consequently, \eqref{eq:PC1} is the trace of the $r$th power of Frobenius on
primitive middle cohomology. This exact complete-intersection
decomposition and trace computation also appear in the proof of
Deligne's Theorem~(8.1), pp.~301--302; the trace formula is
(1.5.1), p.~275 \citep{Deligne1974}.
Deligne's Theorem~(1.6), p.~276, makes the middle characteristic
polynomial integral. Dividing out the middle ambient factor, when
present, preserves integrality.

Katz--Messing's Theorem~1, p.~74 \citep{KatzMessing1974}, identifies
the characteristic polynomials in crystalline and $\ell$-adic
cohomology for smooth projective absolutely irreducible varieties.
The crystalline theory is explicitly included in the statement's
preceding paragraph. Thus the $p$-adic valuations of the primitive
reciprocal roots are its crystalline slopes.

For these smooth complete intersections, the Artin--Mazur comparison
hypotheses were verified above. III, Corollary~3.4, p.~119
\citep{AM}, describes the slopes $a_i/b_i<1$, written
with multiplicity $b_i$, by
\[
 \sum_i(b_i-a_i)\le h^{0,n}=1,
\]
with equality exactly in finite formal height.
Each nonempty summand is a positive integer. In finite height there
is therefore just one slope below one; since its multiplicity is the
height, it is $(h-1)/h$ with multiplicity $h$.
In infinite height the inequality is strict, so there is no slope
below one. The primitive part has the same slopes below one: the
discarded middle ambient class, if any, has slope $n/2\ge1$.

\begin{theorem}[Point-count test on the Hasse-zero locus]
\label{thm:height-two-point-counts}
Assume that the formal height of $X$ is at least two.
Set
\[
 e_2=\frac{\tau_1^2-\tau_2}{2}\in\mathbb Z.
\]
Then
\begin{equation}
 h(X)=2\quad\Longleftrightarrow\quad v_p(e_2)=1,       \label{eq:PC2}
\end{equation}
where $v_p(0)=+\infty$.
For odd $p$, this is equivalent to
\begin{equation}
 \tau_2/p\not\equiv0\pmod p,                         \label{eq:PC3}
\end{equation}
which requires $N_2$ only modulo $p^2$, since
$\sum_{i=0}^{n}p^{2i}\equiv1\pmod{p^2}$.
For $p=2$, the equivalent test is
\begin{equation}
 \frac{\tau_1^2-\tau_2}{4}\not\equiv0\pmod2,          \label{eq:PC4}
\end{equation}
which is determined by $N_1\bmod4$ and $N_2\bmod8$.
\end{theorem}
\begin{proof}
Let $\alpha_1,\ldots,\alpha_B$ be the primitive Frobenius eigenvalues
in an algebraic closure of $\mathbb Q_p$, counted with multiplicity.
Newton's identity gives
\[
 e_2=\sum_{i<j}\alpha_i\alpha_j
     =(\tau_1^2-\tau_2)/2,
\]
and the integral characteristic polynomial proves $e_2\in\mathbb Z$.

If $h=2$, precisely two eigenvalues have valuation $1/2$, while all
remaining eigenvalues have valuation at least one. Their product is
the unique term of valuation one in $e_2$; every other term has
valuation strictly greater than one. Therefore $v_p(e_2)=1$.

If $3\le h<\infty$, every eigenvalue has valuation at least
$1-1/h>1/2$, so every pair product has valuation greater than one.
Integrality implies $v_p(e_2)\ge2$. If $h=\infty$, every valuation is
at least one, which gives the same conclusion.
This proves \eqref{eq:PC2}.

Under the hypothesis $h\ge2$, every eigenvalue has positive valuation.
Therefore the integer $\tau_1$ is divisible by $p$.
For odd $p$,
\[
 e_2\equiv-\tau_2/2\pmod{p^2},
\]
so \eqref{eq:PC2} is equivalent to \eqref{eq:PC3}. The quotient is integral, since
$\tau_2=\tau_1^2-2e_2$ is divisible by $p$.

At $p=2$, $e_2$ is even under the same hypothesis, and
$e_2/2=(\tau_1^2-\tau_2)/4$.
An even $\tau_1$ has square modulo eight determined by its residue
modulo four. Thus \eqref{eq:PC4} uses exactly the stated sufficient precisions.
In terms of counts alone, because $n\ge1$,
\begin{equation}
 \frac{(N_1-3)^2-(-1)^n(N_2-5)}4\pmod2             \label{eq:PC5}
\end{equation}
is the same test: the ambient sums are $3\bmod4$ and $5\bmod8$.
\end{proof}

The height-at-least-two hypothesis is essential to this formulation.
The characteristic-two formula cannot in general be reduced to
point counts modulo $p^2=4$, as the following smooth examples show.

\subsection{An exact characteristic-two precision obstruction}
Over $\mathbb F_2$, in variables $x_1,x_2,x_3,y,z$, define
\[
 F_A=yzA+y^5+x_1^5+x_2^5+x_3^5-x_1z^4,
\]
with either
\[
 A_0=x_1^3,\qquad A_2=x_1^3+x_2^3+x_3^3.
\]
Both equations define smooth quintic threefolds.
Indeed,
\[
 (F_A)_y=zA+y^4,\qquad (F_A)_z=yA,\qquad
 y(F_A)_y-z(F_A)_z=y^5.
\]
At a common zero of all partial derivatives, $y=0$ and the other
partials force $x_2=x_3=0$ and $x_1^4=z^4$.
At a projective point this makes $x_1,z$ nonzero, whereas
$(F_A)_y=z x_1^3\ne0$, a contradiction.

For both forms $a_2=0$. For $a_4$, whose target has exponent three
in every variable, none of the pure fifth powers or the term
$x_1z^4$ can occur. Consequently
\[
 a_4(F_A)=[(x_1x_2x_3)^3]A^3
 =\begin{cases}0,&A=A_0,\\6,&A=A_2.\end{cases}
\]
The logarithm criterion therefore gives height at least three for
$X_{A_0}$ and height exactly two for $X_{A_2}$.

Exact enumeration over $\mathbb F_2$ and
$\mathbb F_4=\mathbb F_2[t]/(t^2+t+1)$ gives
\[
\begin{array}{c|rrrrrr}
 A&N_1&N_2&\tau_1&\tau_2&e_2&h\\\hline
 A_0&15&85&0&0&0&\ge3\\
 A_2&15&89&0&-4&2&2
\end{array}
\]
Projective representatives were enumerated over both fields, and the
identity $C_q=1+(q-1)N_q$ was checked separately on the affine cones.
The reproducible integer and finite-field calculations are included
with the other programs in the workbench.


The examples have the same $N_1$ even as an integer and the same
$N_2\bmod4$, but different height-two status. Hence those coarser
point-count data do not decide height two, even within smooth
anticanonical hypersurfaces. Their $N_2\bmod8$ values differ, as
the proved criterion requires.

\section{Exact height: termination and complexity}
We use division-free arithmetic circuits over $\mathbb F_p$, with the
$d$ variables explicitly listed, for a homogeneous degree-$d$ equation
of a promised smooth hypersurface $X\subset\mathbb P^{d-1}$, $d\ge3$.
The output of the exact-height problem is a binary integer or $\infty$.
The prime is fixed in the modular counting assertions; the polynomial-space
assertion also permits a binary-encoded input prime.  Set $n=d-2$.

\begin{proposition}[A finite stopping bound]
Every finite height of $X$ is at most
\begin{equation}
 B_d=\begin{cases}
 2&d=3,\\
 36&d=4,\\
 \binom{2d-1}{d-1}&d\ge5.
 \end{cases}
\label{eq:HC1}
\end{equation}
\end{proposition}
\begin{proof}
Van der Geer and Katsura (2003, Corollary~2.3, p.~100) prove
$h\le h^{n-1}(\Omega_X^1)+1$ for finite $h$.  Serre duality and
$\omega_X\simeq\mathcal O_X$ identify this dimension with $h^1(T_X)$.
For $n\ge3$, the restricted Euler sequence and hypersurface line-bundle
cohomology give $H^1(T_{\mathbb P^{d-1}}|_X)=0$: indeed
$H^1(\mathcal O_X(1))=H^2(\mathcal O_X)=0$.  The normal sequence
then gives
\[
 h^1(T_X)\le h^0(\mathcal O_X(d))
 =\binom{2d-1}{d-1}-1.
\]
For $d=4$ the same Euler sequence gives
$h^1(T_{\mathbb P^3}|_X)=h^2(\mathcal O_X)=1$, so that
$h^1(T_X)\le34+1=35$.  For $d=3$,
$h^0(\Omega_X^1)=1$.  This proves \eqref{eq:HC1}, without assuming that
infinitesimal automorphisms vanish or that $p$ does not divide $d$.
\end{proof}

\begin{proposition}[Exact modular coefficient algorithm]
Let $\widetilde F$ be any homogeneous integral lift and put
\[
 a_{p^j}=[(x_1\cdots x_d)^{p^j-1}]\widetilde F^{p^j-1}.
\]
Compute $a_{p^j}\bmod p^j$ successively for $1\le j\le B_d$.
The first nonzero residue has index $h(X)$; if all residues vanish,
$h(X)=\infty$.
\end{proposition}
\begin{proof}
Lemma~\ref{lem:logarithm-height}
says that conditional on $h\ge j$ one has
$v_p(a_{p^j})\ge j-1$, with equality exactly when $h=j$.
Thus the first nonzero residue has precisely the claimed meaning.
If all residues vanish, a finite height would exceed \eqref{eq:HC1}, which is
impossible.  Only precision $p^j$ is required at the $j$th step.

For a circuit implementation, first form the homogeneous components
through degree $d$ and keep the degree-$d$ output. Choose integer
representatives of its scalar literals once and fix this homogeneous
integral circuit lift for all steps. A literal lift of the original
circuit need not be homogeneous, because cancellations in unwanted
degrees may have occurred only modulo $p$.
Evaluate the resulting circuit in the finite ring
\[
 (\mathbb Z/p^j\mathbb Z)[x_1,\ldots,x_d]/
 (x_1^{p^j},\ldots,x_d^{p^j}),
\]
raise its value to $p^j-1$ by repeated squaring, and read the diagonal
coefficient. Nonnegative exponents make truncation exact.
\end{proof}

\subsection{The quasi-\texorpdfstring{$F$}{F}-splitting matrix and complexity}
\label{subsec:quasi-height-complexity}

The full-height iteration used here is an existing algorithm.
Kawakami, Takamatsu, and Yoshikawa give its exact criterion
\citep[Theorem~C, version~3, p.~4; Theorem~5.8, pp.~37--39]{KTY2025};
the criterion was already present in version~2 of their preprint, p.~5.
Batubara, Garzella, and Pan identify the fixed finite-dimensional operator
and implement its matrix iteration
\citep[Section~3.4 and Algorithm~4, p.~7]{BGP2025}.
Yobuko identifies quasi-$F$-split height with Artin--Mazur height for
Calabi--Yau varieties
\citep[Theorem~4.5; author version, pp.~9--10]{Yobuko2019}.
We analyze this algorithm under the circuit input convention of this paper.

Put $S=\mathbb F_p[x_1,\ldots,x_d]$ and write
\[
 D=d(p-1),\qquad g=F^{p-1},\qquad
 V=S_D,\qquad a=(p-1,\ldots,p-1),
\]
\[
 W=\dim_{\mathbb F_p}V=\binom{dp-1}{d-1}.
\]
Thus $V=S_D$ is the space of homogeneous degree-$D$ polynomials.
Choose a homogeneous integral coefficient lift
$G=\sum_{|\beta|=D}\widetilde c_\beta x^\beta$ of $g$, and define
\begin{equation}\label{eq:quasi-delta}
 \Delta=\left(
 \frac{G^p-\sum_{|\beta|=D}\widetilde c_\beta^{\,p}x^{p\beta}}p
 \right)\bmod p.
\end{equation}
This is the monomial representative $\Delta_1(g)$ of the cited criterion;
see also \citet[Proposition~3.2, Remark~3.4 and Algorithm~1, p.~5]{BGP2025}.
The sum is over distinct collected monomials.
The numerator is divisible by $p$ coefficientwise. Moreover, replacing
$G$ by $G+pH$, or a scalar lift by $\widetilde c_\beta+pb$, changes the
corresponding $p$th power by a multiple of $p^2$. Thus
\eqref{eq:quasi-delta} is independent of the chosen integer coefficient
lifts. This includes $p=2$ and requires no Teichm\"uller representatives.
Computationally, form the numerator modulo $p^2$, choose each residue in
$[0,p^2)$, and divide it by $p$; the resulting residue in $[0,p)$ is exact.

Define the $\mathbb F_p$-linear map $U$ on monomials by
\[
 U(x^\gamma)=
 \begin{cases}
 x^{(\gamma-a)/p},&\gamma\equiv a\pmod p,\\
 0,&\text{otherwise}.
 \end{cases}
\]
Since $\deg\Delta=pD$, the map
\[
 M:V\longrightarrow V,\qquad v\longmapsto U(\Delta v)
\]
preserves $V$: its retained monomials have degree
$(pD+D-d(p-1))/p=D$. In the monomial basis,
\begin{equation}\label{eq:quasi-matrix-entry}
 M_{\alpha,\beta}=[x^{p\alpha+a-\beta}]\Delta,
 \qquad |\alpha|=|\beta|=D,
\end{equation}
where a negative exponent gives zero. Let $\ell(v)=[x^a]v$.
The only degree-$D$ monomial outside $(x_1^p,\ldots,x_d^p)$ is $x^a$.
Consequently the cited height criterion gives
\begin{equation}\label{eq:quasi-height-iteration}
 h(X)=\inf\{j\ge1:\ell(M^{j-1}g)\ne0\},\qquad \inf\varnothing=\infty.
\end{equation}

For the field hypotheses, apply Kawakami, Takamatsu, and Yoshikawa's
projective criterion and Yobuko's
comparison after extension to $\overline{\mathbb F}_p$. The Cartier map
over a perfect field takes $p$th roots of coefficients, but every
coefficient of this recurrence lies in $\mathbb F_p$, so the recurrence
is exactly the linear one displayed above. Formal-group height is
unchanged by extending the base field, since nonvanishing of the first
nonzero coefficient of the $p$-series is unchanged. The adjunction and
cohomology calculations already made verify the Calabi--Yau hypotheses.
This applies also when $d=3$.

\begin{lemma}[A matrix stopping bound]
\label{lem:quasi-matrix-stopping}
In \eqref{eq:quasi-height-iteration}, a finite height is at most $W$.
Thus testing $j=1,\ldots,W$ detects infinite height as well. For the
smooth inputs here one may use the smaller cutoff $\min(W,B_d)$.
\end{lemma}
\begin{proof}
The sequence $y_r=\ell(M^rg)$ satisfies a linear recurrence of order
at most $W$ by Cayley--Hamilton. If its first $W$ terms vanish, the
recurrence forces every subsequent term to vanish. Otherwise the first
nonzero term has index at most $W-1$, which proves the assertion by
\eqref{eq:quasi-height-iteration}. This argument requires no computation
of the characteristic polynomial. The alternative bound $B_d$ is
\eqref{eq:HC1}.
\end{proof}

Write $\mathrm{FPSPACE}$ for function problems computable deterministically
with a number of work bits polynomial in the input length.

\begin{theorem}[Resource bounds for exact height]
\label{thm:quasi-height-resource-bounds}
For a binary-encoded prime $p$ and a division-free circuit of size $s$
in $d\ge3$ explicitly listed variables, promised to define a smooth
degree-$d$ hypersurface in $\mathbb P^{d-1}$, exact Artin--Mazur height,
including $\infty$, belongs to $\mathrm{FPSPACE}$.
There is also a deterministic dense implementation with bit complexity
\[
 2^{O(d\log p)}\operatorname{poly}(s,d,\log p).
\]
The polynomial-space and dense-time bounds refer to separate
implementations of the same height criterion.
\end{theorem}
\begin{proof}
Let $L$ denote input length and set $E=pD=dp(p-1)$. All scalar
calculations below take place modulo $p$ or $p^2$, so each scalar has
$O(\log p)$ bits. A monomial of total degree at most $E$ has an exponent
vector using $O(d\log(E+1))$ bits. Furthermore,
\[
 \log W=O(d(1+\log p)),\qquad
 \log(E+1)=O(\log d+\log p).
\]
These lengths are polynomial in $L$, because $d$ is explicitly bounded
by input length.

\paragraph{Coefficient access.}
For a degree-$d$ output coefficient of $F$, recurse through its arithmetic
circuit. At a product gate, enumerate all coordinatewise splits of the
target exponent vector, obtain the two child coefficients, and accumulate
their product modulo $p$. At a sum gate, add the child coefficients.
Constants and variables are immediate. A recursion path has at most $s$
gates, and each frame stores polynomial-length exponent vectors, gate
addresses, a loop counter and an $O(\log p)$-bit accumulator. Intermediate
gates need not be homogeneous. Equivalently, this computation evaluates
the circuit modulo all monomials of total degree greater than $d$, which
retains the required output exactly. The degree-$d$ output is selected
using the promise; an output coefficient of any other total degree is zero.

Coefficients of a power $A^r$ are obtained by the same convolution
recursion applied to
\[
 A^r=A^{\lfloor r/2\rfloor}A^{\lceil r/2\rceil},
\]
with base cases $r=0,1$. The power recursion has depth
$O(\log(r+1))$. First compute coefficients of $g=F^{p-1}$ modulo $p$.
Treat their representatives in $[0,p)$ as the coefficients of a virtual
homogeneous lift $G$, and compute coefficients of $G^p$ modulo $p^2$.
The two powering recursions add only $O(\log p)$ depth to the original
circuit recursion. In the pure-power sum in \eqref{eq:quasi-delta}, a
requested exponent has at most one contributing term, determined by
coordinatewise divisibility by $p$. Scalar powering modulo $p^2$ and
the exact scalar division by $p$ complete the coefficient calculation.
Thus every coefficient of $\Delta$, and hence every matrix entry in
\eqref{eq:quasi-matrix-entry}, is accessible in polynomial space.

\paragraph{Matrix powers.}
For $r\ge2$, compute one entry recursively from
\[
 (M^r)_{\alpha,\beta}
   =\sum_{|\gamma|=D}
      (M^{\lfloor r/2\rfloor})_{\alpha,\gamma}
      (M^{\lceil r/2\rceil})_{\gamma,\beta}\pmod p.
\]
The cases $r=0,1$ use the identity and the entry oracle just constructed.
For $r<W$, the matrix recursion has depth $O(\log W)$.
Every frame stores two index vectors, the current summation index, a
binary exponent, and one field accumulator. All have polynomial length.
An entry-oracle call adds only the polynomial stack used for coefficients.
Finally, compute
\[
 \ell(M^{j-1}g)=
 \sum_{|\beta|=D}(M^{j-1})_{a,\beta}g_\beta\pmod p
\]
and scan $j=1,\ldots,W$ using a binary counter. Return the first nonzero
test, or $\infty$ if all vanish. Lemma~\ref{lem:quasi-matrix-stopping}
proves correctness. The outer scan does not create a recursion chain of
length $W$; every requested power uses divide and conquer. The complete
stack, counters, exponents and scalar accumulators use $L^{O(1)}$ space.

\paragraph{A separate dense time bound.}
The number of monomials of total degree at most $E$ is
\[
 V_E=\binom{E+d}{d}
 \le \bigl(e[p(p-1)+1]\bigr)^d
 \le (ep^2)^d,
 \qquad W\le\binom{dp}{d}\le(ep)^d.
\]
Expand the original circuit modulo $p$ through total degree $d$,
using at most $\binom{2d}{d}\le4^d$ positions. Dense binary powering
computes $g$ and $G^p$ in $O(\log p)$ polynomial-product stages.
Naive coefficient convolution uses at most $V_E^2$ scalar products
per stage. Form $\Delta$ by \eqref{eq:quasi-delta} and construct the
$W^2$ matrix entries by coefficient lookup. At most $W-1$ matrix-vector
products then test the whole sequence, using $O(W^3)$ field operations.
Degree indexing and scalar arithmetic contribute a fixed polynomial in
$d,\log p$; even sequential array-access implementations contribute
only a fixed polynomial in $V_E,W$. Thus the total bit complexity is
a fixed polynomial in $s,d,\log p,V_E,W$, giving
$2^{O(d\log p)}\operatorname{poly}(s,d,\log p)$ uniformly in all
parameters. This dense implementation uses exponentially large arrays,
whereas the preceding recursion recomputes entries to obtain the
polynomial-space bound.
\end{proof}



\subsection{Uniform characteristic and exact counting}
Varying the characteristic gives a stronger lower bound than fixing one modular counting class. We state the encoding explicitly, since this distinction is essential.

The symbols $\mathrm P$ and $\mathrm{FP}$ denote deterministic
polynomial-time decision and function computation, respectively.
A superscript denotes access to the indicated oracle; in particular,
a $\#\mathrm P$ oracle returns an exact integer count.

\begin{theorem}[Uniform height two and exact counting]\label{thm:uniform-counting}
Let $\mathcal H_2$ be the promise decision problem whose input is a prime $p$ in unary and a width-three formula $\Phi$, promised to have $p\mid\#\operatorname{SAT}(\Phi)$, and whose answer is whether the counting hypersurface $X_\Phi$ has height two. Then
\[
 \mathrm{FP}^{\mathcal H_2}=\mathrm{FP}^{\#\mathrm P},
\]
where oracle calls to the promise problem must satisfy its promise.
The reduction from exact counting uses polynomially many nonadaptive calls. Consequently the full exact-height function for the circuit presentations of this paper is $\#\mathrm P$-hard under polynomial-time nonadaptive oracle reductions, even when every query has height at least two and its characteristic is polynomially bounded in the size of the original counting input.
\end{theorem}
\begin{proof}
For the upper inclusion, the ordered local lists defining $T_{p+1}(\Phi)$ have length polynomial in the unary input length and are checked by the displayed marginal equations. Hence this integer is a $\#\mathrm P$ count uniformly in $p$. On the promise, the oracle answer is obtained from its exact value by reducing modulo $p^2$ and testing whether the residue is nonzero. The height-two criterion proves correctness. An exact-counting oracle therefore simulates each promised $\mathcal H_2$ call in polynomial time.

For the other inclusion let $\Phi$ have $N\ge1$ explicitly declared variables, and put $S=\#\operatorname{SAT}(\Phi)$, so $0\le S\le2^N$. For any polynomially bounded positive prime $p$ and any $a\in\{0,\ldots,p-1\}$, construct a width-three formula $\Phi_a$ with exactly $S+a$ satisfying assignments as follows. Introduce a branch bit $b$ and $k=\lceil\log_2 p\rceil$ bits $z$, and use the Boolean predicate
\[
 (\neg b\wedge\Phi(x)\wedge z=0)
 \ \vee\
 (b\wedge x=0\wedge z<a).
\]
The two branches have respectively $S$ and $a$ assignments. The comparison $z<a$ has a Boolean circuit of size polynomial in $k$; the whole predicate therefore has polynomial size. Give each circuit gate a new variable and impose its truth table by clauses of width at most three, then force the output to one. For example, $v=u\wedge w$ is enforced by
\[
 (\neg v\vee u)\wedge(\neg v\vee w)
       \wedge(v\vee\neg u\vee\neg w),
\]
and a negation uses two binary clauses. OR gates are treated by the dual three clauses. Each input assignment has exactly one extension to the gate variables, so this transformation preserves the satisfying count exactly. The case $a=0$ uses the always-false second branch. Declared variables and auxiliary variables remain part of the counted assignment.

For each prime use the height-two gadget $\Psi_p$ already proved above: for odd $p$ take $\Theta_{p-1,p-2}$, and for $p=2$ take the displayed three-clause characteristic-two gadget. It has $O(p^2)$ clauses, $p\mid\#\operatorname{SAT}(\Psi_p)$, and
\[
 v_p\bigl(T_{p+1}(\Psi_p)\bigr)=1.
\]
Conjoin $\Phi_a$ with $\Psi_p$ on disjoint variable sets. Every resulting query satisfies the Hasse-zero promise. Multiplicativity and $T_{p+1}(\Phi_a)\equiv(S+a)^2\pmod p$ give
\[
 h(X_{\Phi_a\wedge\Psi_p})=2
 \quad\Longleftrightarrow\quad S+a\not\equiv0\pmod p.
\]
Exactly one of the $p$ answers is therefore negative, identifying $S\bmod p$. These $p$ queries are fixed before any answer is read.

We give an elementary polynomial bound for the primes required. Put $M=4(N+1)^2$, let $P_M$ be the product of the primes at most $M$, and let $L_M=\operatorname{lcm}(1,\ldots,M)$. Every binomial coefficient $\binom Mj$ divides $L_M$: for each prime, its factorial-valuation formula is a sum of at most $\lfloor\log_p M\rfloor$ terms, each zero or one. Hence
\[
 L_M\ge\max_j\binom Mj\ge\frac{2^M}{M+1},
 \qquad L_M\le P_M^{\lfloor\log_2M\rfloor}.
\]
For $N\ge1$, $M\ge16$, $\log_2(M+1)\le M/2$, and $\log_2M\le2(N+1)$. These inequalities imply $P_M\ge2^{N+1}>2^N$. Trial division or a sieve lists all primes up to $M$ in polynomial time. Asking all shifts for all these primes takes at most $M^2$ queries. Chinese remaindering then recovers the unique $S\in[0,2^N]$ from its residues. The case $N=0$ can be evaluated directly.

Every query has polynomial length. The gadget has $O(p^2)$ incidences, while the sum construction has polynomial size in the original formula length and $\log p$. In the hypersurface construction the least power satisfies $Q\le p(u+2)$, so its explicitly listed $Q+1$ variables and its arithmetic circuit still have polynomial size because $p\le M$. Thus the reduction also supplies valid polynomial-size circuit inputs to the exact-height problem.

Finally, any $\#\mathrm P$ function counts the accepted strings of a polynomial-time verifier of polynomially bounded length. Unrolling that verifier into a Boolean circuit and applying the unique-extension gate clauses just used gives a parsimonious width-three formula. The reduction consequently applies to all $\#\mathrm P$ functions. Replacing each decision query by an exact-height query and testing whether the answer is two proves the final assertion. The value $\infty$ is simply one of the negative answers. The finite-height bound gives $O(d)$ bits for every finite answer, so reading a height answer also costs polynomial time.
\end{proof}
The prime varies in this theorem. It supplies no $\#\mathrm P$-hardness statement at one fixed characteristic, and its exact-counting equivalence concerns the uniform height-two oracle. The complexity of the unrestricted full-height function still requires an upper bound matching its lower bound.


\subsection{Height supplied in unary}
The matrix criterion also permits modular counting when the requested
height grows with the input. The relevant witness is a path in the
matrix, rather than a coefficient of degree exponential in that height.

\begin{theorem}[Unary height decisions]\label{thm:unary-height}
Fix a prime $p$. On the smooth homogeneous circuit inputs specified
above, with a positive integer $r$ supplied in unary, deciding
$h(X)=r$ is $\mathrm{Mod}_p\mathrm P$-complete under polynomial-time
many-one reductions. In particular, membership holds uniformly in $r$;
the completeness assertion for this variable-$r$ problem already holds
on its restriction $r=2$.

Let $H_{\mathrm{unary}}$ denote the same decision problem with both
$p$ and $r$ supplied in unary, and with $p$ promised prime. Then
$H_{\mathrm{unary}}\in\mathrm P^{\#\mathrm P}$, using one exact
counting query followed by a remainder computation. With
promise-respecting oracle calls,
\[
 \mathrm{FP}^{H_{\mathrm{unary}}}=\mathrm{FP}^{\#\mathrm P}.
\]
\end{theorem}
\begin{proof}
We first construct nonnegative counts for the coefficients entering
the operator in \eqref{eq:quasi-delta}. Homogenize the input circuit
through degree $d$ and keep its degree-$d$ output. Evaluate each
degree-zero subcircuit in $\mathbb F_p$ and replace it by its residue
in $\{0,\ldots,p-1\}$. Rewrite subtraction as addition after
multiplication by $p-1$. The resulting arithmetic circuit uses only
addition, multiplication and nonnegative integer scalar literals;
it represents an integral homogeneous lift $\widetilde F$ of $F$.
Repeated squaring constructs a circuit for
\[
 G=\widetilde F^{p-1}=\sum_{|\beta|=D}c_\beta x^\beta,
 \qquad D=d(p-1),\qquad c_\beta\ge0.
\]
Its size is polynomial in the input length, and $D$ is polynomial
when $p$ is fixed or unary.

A parse witness at an addition chooses one child, at a product takes
both children, and at a scalar $b$ chooses one of $b$ labels. The
exponent vector records the variable leaves. Thus the number of
witnesses with vector $\beta$ is exactly $c_\beta$. These are
polynomial-length witnesses even though the circuit may share
subcomputations. After degree-zero gates have been collapsed, every
positive-degree leaf contributes at least one to total degree. There
are at most $D$ such leaves and at most $D-1$ product nodes having
two positive-degree children. Between those branching nodes a chain
has length at most the circuit size; scalar children add at most
one leaf per node. Hence a parse has polynomially many nodes,
bounded by the circuit size times $O(D)$. Choose a fixed-length
encoding with unique padding, so each parse contributes exactly one
accepting witness and the witnesses have a canonical lexical order.
This proves that $c_\beta$ is a $\#\mathrm P$ count, uniformly in
the indicated input parameters.

Consider the integral polynomial
\[
 Q=\frac{G^p-\sum_{|\beta|=D}c_\beta^{\,p}x^{p\beta}}p.
\]
The sum is over distinct collected monomials. For an exponent vector
$\gamma$, the numerator coefficient counts ordered $p$-tuples of
parse witnesses whose vectors sum to $\gamma$, excluding tuples
whose exponent vectors are all equal. The excluded number is
precisely $\sum_{p\beta=\gamma}c_\beta^p$: it includes all parses
of the same monomial, not only tuples of identical parses.

Cyclic rotation acts freely on the remaining tuples. Indeed, since
$p$ is prime, a tuple fixed by a nonidentity rotation has all its
parse witnesses equal, and therefore has all exponent vectors equal.
Consequently $[x^\gamma]Q$ is the number of these free orbits.
A verifier guesses the $p$ parses, checks their exponent sum and
that their vectors are not all equal, and accepts exactly the
lexically least cyclic rotation. Both the witness and the comparison
with all $p$ rotations have polynomial length and running time for
fixed or unary $p$. Every coefficient of $Q$ is thus itself a
$\#\mathrm P$ count. As a small example, if $G=cx^\beta+ex^\eta$
with $\beta\ne\eta$ and $p=2$, then
$Q=cex^{\beta+\eta}$: the two possible orders of one parse of
each monomial form one orbit.

By coefficient-lift independence in \eqref{eq:quasi-delta},
$Q\bmod p=\Delta$. In particular a coefficient of $Q$ gives a
nonnegative counting lift of each matrix entry
\[
 M_{\alpha\beta}=[x^{p\alpha+a-\beta}]\Delta.
\]
A negative component of $p\alpha+a-\beta$ gives zero witnesses.
For $j\ge1$, expand the scalar
\[
 y_j=\ell(M^{j-1}g)
\]
as a sum over paths of $j-1$ matrix transitions ending at the
monomial $x^a$. A path witness specifies all its degree-$D$
monomial indices, a parse witnessing its initial coefficient of
$G$, and an orbit witness for each transition. Each index has
$O(d\log(D+1))$ bits. Therefore, for unary $j$, this is a
polynomial-length witness, with a polynomial-time verifier.
Let $N_j$ be its number of accepting witnesses. Summing over paths
and multiplying the choices along each path proves
\[
 N_j\equiv y_j\pmod p.
\]
This construction includes $j=1$, when the path has no transitions
and the count is the coefficient of $x^a$ in $G$.

Nonnegative counting functions are closed under sums and under
polynomially many products: a witness selects a summand or supplies
one witness for each factor, respectively. Thus
\begin{equation}\label{eq:unary-count}
 Z_r=N_r^{p-1}\prod_{j=1}^{r-1}
       \bigl(1+(p-1)N_j^{p-1}\bigr)
\end{equation}
is a $\#\mathrm P$ count when $p$ is fixed or unary and $r$ is
unary. The sum of the path lengths in this product is $O(r^2)$,
multiplied by polynomial factors in the unary prime and circuit
length. Fermat's identity in $\mathbb F_p$ gives
\[
 Z_r\not\equiv0\pmod p
 \quad\Longleftrightarrow\quad
 y_1=\cdots=y_{r-1}=0\ \hbox{ and }\ y_r\ne0.
\]
The matrix criterion identifies this condition with $h(X)=r$.
For fixed $p$ this proves $\mathrm{Mod}_p\mathrm P$ membership;
the height-two reduction proves completeness of the variable-$r$
problem. If both $p$ and $r$ are unary, an exact counting oracle
computes $Z_r$, whose binary length is polynomial in the witness
bound, and the remainder modulo $p$ decides the problem. This
proves the claimed $\mathrm P^{\#\mathrm P}$ upper bound.
Conversely, every query to the promised uniform height-two oracle
in Theorem~\ref{thm:uniform-counting} is a valid query to
$H_{\mathrm{unary}}$ with $r=2$. That theorem recovers every
$\#\mathrm P$ function. Composing these two simulations proves
the stated equality of function-oracle classes.
\end{proof}

\begin{corollary}[Unary thresholds]\label{cor:unary-threshold}
For a fixed prime $p$ and unary $R$, deciding $h(X)\le R$ belongs
to $\mathrm{Mod}_p\mathrm P$. If a specified polynomial $q$ is
fixed and the input additionally promises that every finite height
is at most $q(L)$, where $L$ is circuit input length, then the
exact-height function, including $\infty$, belongs to
$\mathrm{FP}^{\mathrm{Mod}_p\mathrm P}$ on that promise.
\end{corollary}
\begin{proof}
Over $\mathbb F_p$, the expression
$1-\prod_{j=1}^R(1-y_j^{p-1})$ is nonzero exactly when some
$y_j$ is nonzero. Replace each negative scalar by its residue
$p-1$ and use the counts just constructed; the result is a
$\#\mathrm P$ count modulo $p$ with polynomial witness length.
For the second assertion, query the exact-height predicate for
$r=1,\ldots,q(L)$. Return its first positive answer or $\infty$
if none occurs. The additional promise makes the latter conclusion
valid. No polynomial bound on finite heights is assumed for
general circuit inputs, and the argument does not allow a
binary-encoded height parameter.
\end{proof}


\paragraph{What the lower bounds establish.}
The counting construction makes height one
$\mathrm{Mod}_p\mathrm P$-complete for every fixed prime.  Combined
with the fixed-gadget reduction above, the preceding
membership theorem makes height two $\mathrm{Mod}_p\mathrm P$-complete
as well.  At each fixed characteristic, the full exact-height function is therefore
$\mathrm{Mod}_p\mathrm P$-hard. With the characteristic as part of the
input, Theorem~\ref{thm:uniform-counting} strengthens the lower bound
to $\#\mathrm P$-hardness under nonadaptive polynomial-time oracle
reductions. The full function lies in $\mathrm{FPSPACE}$, and hence in
$\mathrm{FEXPTIME}$, the function class with deterministic running time
$2^{\operatorname{poly}(L)}$ for input length $L$.
These are distinct bounds; they do not constitute a complete complexity
classification of the full function. A single height-one zero test
does not determine the full residue of a counting function.
Theorem~\ref{thm:unary-height} supplies uniform modular counting for
unary height parameters, but does not decide infinite height in that
class or extend the witness bound to binary height parameters.

\section{The Fermat character and the first variation}
\label{subsec:fermat-first-variation}

Throughout this section, characters refer to coordinate pullback.
Let $k=\overline{\mathbb F}_2$, let $r\geq3$ be odd, and put
\[
 H=z_1^r+\cdots+z_r^r,\qquad X_0=V(H)\subset\mathbb P_k^{r-1}.
\]
The group $D=\mu_r^r$ acts by $z_i\mapsto\lambda_i z_i$.
Its diagonal subgroup acts trivially on projective space.

\begin{proposition}[The top coherent character]
Put $E_0=H^{r-2}(X_0,\mathcal O_{X_0})$.
The character of $D$ on this line is
\[
 \chi_{\rm top}(\lambda_1,\ldots,\lambda_r)
       =(\lambda_1\cdots\lambda_r)^{-1}.
\]
This character descends to $D/\mu_r$.  Absolute $p$-power
Frobenius carries a $\chi$-eigenvector to a $\chi^p$-eigenvector.
In particular, for $p=2$ the absolute Frobenius on $E_0$ is zero.
\end{proposition}
\begin{proof}
Since $H$ is $D$-invariant, the hypersurface sequence
\[
 0\longrightarrow\mathcal O_{\mathbb P^{r-1}}(-r)
   \xrightarrow{\ H\ }\mathcal O_{\mathbb P^{r-1}}
   \longrightarrow\mathcal O_{X_0}\longrightarrow0
\]
is equivariant.  Its connecting map identifies $E_0$ with
$H^{r-1}(\mathbb P^{r-1},\mathcal O(-r))$, whose standard
\v Cech generator is $e=(z_1\cdots z_r)^{-1}$.
Pullback multiplies this generator by $(\prod_i\lambda_i)^{-1}$.
On a diagonal element $(\lambda,\ldots,\lambda)$ this scalar is
$\lambda^{-r}=1$, which proves descent.

Naturality of absolute Frobenius and its semilinearity give
$g^*F(e)=F(g^*e)=\chi(g)^pF(e)$.  A nonzero Frobenius image
in the same one-dimensional space would therefore require
$\chi^p=\chi$.  The displayed character has order $r$, whereas
for $p=2$ this requirement would make it trivial.  As $r\geq3$,
the Frobenius image is zero.
\end{proof}

Let now $F$ be any degree-$r$ form over $\mathbb F_2$, let
$c=[z_1\cdots z_r]F\in\mathbb F_2$, and consider the smooth open
part of the pencil
\[
 X_u=V(H+uF),\qquad
 U=\operatorname{Spec}\mathbb F_2[u,\Delta(u)^{-1}],\quad
 \Delta(0)=1.
\]
The earlier parity-counting pencil is a special case \citep{ParityCompanion}.
Here the existence of $U$ can be seen directly. Put $d_0=r(r-2)+1$.
Every degree-$d_0$ monomial is divisible by some $z_i^{r-1}$.
For each such monomial, choose one such index and multiply
$\partial_i(H+uF)$ by the complementary monomial. In the basis of
all degree-$d_0$ monomials, the resulting square matrix is $I+uB$,
because $r=1$ in characteristic two. Its determinant $\Delta$ has
constant term one. After inverting $\Delta$, every degree-$d_0$
monomial belongs to the partial-derivative ideal, so the derivatives
have no common projective zero. This proves smoothness over $U$.
The hypersurface sequence over $U$ and its Laurent monomial
calculation give a canonical frame $e$ of
$E=R^{r-2}\pi_*\mathcal O_{X_U}$.

\begin{proposition}[The precise first-order statement]
\label{prop:correct-first-variation}
The linearized coherent Frobenius is the bundle map
\[
 \varphi:F_U^*E\longrightarrow E,\qquad
 \varphi(1\otimes e)=uc\,e.
\]
Its derivative at $u=0$ in this frame and parameter is $c$;
the expression $uc$ is its first-order term.  Intrinsically its
first variation belongs to
\[
 T_0^*U\otimes\operatorname{Hom}_k(F_k^*E_0,E_0).
\]
It is not the differential of the absolute Frobenius morphism
of the base, and it is not an ordinary endomorphism of $E_0$.
\end{proposition}
\begin{proof}
The coherent Frobenius formula identifies its matrix coefficient
with $[z_1\cdots z_r](H+uF)=uc$.  This formula follows directly
by applying Frobenius to the hypersurface sequence and the displayed
\v Cech generator; in the smooth setting it is also Katz's
Algorithm~2.3.7.14, pp.~35--36, and Special Case~2.3.7.17,
p.~36 \citep{Katz1972}.  The coefficient of the product monomial in
$H$ is zero, so $\varphi_0=0$.  The first jet of a bundle map
vanishing at the base point is the displayed tensor, proving
the intrinsic assertion.

For clarity, under a frame change $e'=a(u)e$, with $a(0)\ne0$,
the coefficient changes to $a(u)^{p-1}h(u)$, where $p=2$ and
$h(u)=uc$.  Since $h(0)=0$, its derivative changes by
$a(0)^{p-1}$, exactly the change of basis in the indicated
Frobenius-twisted Hom space.  A change of parameter is accounted
for by the cotangent factor.  On the dual numbers
$k[\epsilon]/(\epsilon^2)$, the semilinear map is explicitly
\[
  F((a+b\epsilon)e)=a^2\epsilon c\,e.
\]
By contrast, the differential of absolute base Frobenius
$u\mapsto u^2$ is zero.  These are different operations.
\end{proof}

\paragraph{Equivariance and the limit of the character interpretation.}
For arbitrary $F$, the full diagonal group does not preserve
the generic fiber of the pencil: it sends the direction $F$ to
$g^*F=F(\lambda_1z_1,\ldots,\lambda_rz_r)$.
There is consequently no full-$D$ character decomposition on
the generic fiber defined by this action.  What is well-defined
at the Fermat point is the twisted first variation of
Proposition~\ref{prop:correct-first-variation}.
Its covariance is explicit:
\[
 c(g^*F)=(\lambda_1\cdots\lambda_r)c(F)
             =\chi_{\rm top}(g)^{-1}c(F).
\]
Indeed, among degree-$r$ monomials, the only one of full-$D$
weight $\prod_i\lambda_i$ is $z_1\cdots z_r$: every exponent
must be congruent to one modulo $r$, and their sum is $r$.
Equivalently, the coordinate change $X_{g^*F}\to X_F$ pulls
the canonical frame back to $\chi_{\rm top}(g)e$; compatibility
with the semilinear map gives
$h_{g^*F}=\chi_{\rm top}(g)^{1-p}h_F$.
For $p=2$, the Hom space
$\operatorname{Hom}(F_k^*E_0,E_0)$ has character
$\chi_{\rm top}^{-1}$, as required.

The incidence form has a useful smaller symmetry.
Write it as $F=z_0P$, with $P$ homogeneous of degree two in
each of its $M$ pairs of incidence variables, so $r=2M+1$.
Scale $z_0$ by $\lambda_0$ and both coordinates of pair $e$
by $\lambda_e$.  Then
\[
 g^*F=\rho(g)F,\qquad
 \rho(g)=\lambda_0\prod_{e=1}^M\lambda_e^2
         =\chi_{\rm top}(g)^{-1}.
\]
Together with $u\mapsto\rho(g)^{-1}u$, this defines an action
on the total pencil.  Its subgroup $\ker\rho$ acts fiberwise
and has trivial character on the coherent line.  Thus this
symmetry gives an equivariant family with a nontrivial action
on its parameter, not a full diagonal character action on
every fiber.  The statement concerns the usual incidence form;
the general first-variation proposition also applies to any exceptional cubic chosen in a counting construction.

In particular, the literal assertion that ``$uc$ is the derivative
on the $\chi_{\rm top}$-line'' needs two corrections: the derivative
is $c$, and its source is the Frobenius twist of that line.
Identifying this coherent first variation with a deformation of
a fixed rank-one Jacobi-sum motive would require an additional
motivic construction; the character calculation alone supplies
no such identification.

\subsection{The rank parameter depends on the input model}
\label{subsec:rank-model}

A rank bound does not specify how Frobenius is presented.
If an $R\times R$ Frobenius matrix is given modulo $p$, its
trace requires $R$ diagonal entries and $R-1$ field additions.
By contrast, a succinct geometric formula need not supply a
basis or the matrix.  The rank-one coherent pencil above is
an example of this distinction, but coherent rank is not the
rank of the full pure motive.  The following gives a separate
obstruction for genuine Artin motives, with its syntax stated
explicitly.

Here $\mathrm{NP}$ is nondeterministic polynomial time, and
$\mathrm{RP}$ is randomized polynomial time with one-sided error:
inputs outside the language are always rejected, while inputs in the
language are accepted with probability at least one half.

\begin{proposition}[A promised bounded-rank geometric model]
Fix a finite field $\mathbb F_q$ of characteristic $p$.
Suppose a deterministic algorithm, on polynomial presentations of finite
\'etale schemes $Y/\mathbb F_q$, correctly computes the Frobenius
trace modulo $p$ whenever the rank of $H^0(Y)$ is at most one,
and runs in time polynomial in the presentation size on that
promised class.  Then $\mathrm{NP}=\mathrm{RP}$.
The rank here is a promised upper bound, not an additional
input revealing the exact number of points.
\end{proposition}
\begin{proof}
For a Boolean formula $\Phi$, use Boolean input variables and
one uniquely determined Boolean wire variable for each gate.
Impose $x_i^2-x_i=0$ on inputs and the usual equations
$t-ab=0$, $t-(a+b-ab)=0$, and $t-(1-a)=0$ for AND, OR, and
NOT gates, followed by the equation that the output equals one.
This is a polynomial-size, degree-at-most-two presentation.
Eliminating wire variables identifies its coordinate algebra
with the quotient of
\[
 \mathbb F_q[x_1,\ldots,x_n]/(x_i^2-x_i)_i
       \simeq\prod_{a\in\{0,1\}^n}\mathbb F_q
\]
which retains exactly the satisfying factors.  Hence $Y_\Phi$
is a finite \'etale scheme consisting of $S_\Phi$ rational points.
Its full $H^0$ has rank $S_\Phi$ and Frobenius acts as the
identity, so its trace is $S_\Phi$.

On the promise $S_\Phi\in\{0,1\}$, the proposed algorithm
distinguishes the two cases.  Valiant--Vazirani isolation
therefore implies $\mathrm{NP}=\mathrm{RP}$ (Valiant \& Vazirani,
1986, Theorem~1.1 and Corollary~1.2, printed p.~87).
One may clock the algorithm at its polynomial bound on every
input.  An unsatisfiable input stays unsatisfiable under every
isolation restriction; a satisfiable input becomes uniquely
satisfiable with inverse-polynomial probability.  Arbitrary
answers outside the promise do not create a false positive
from an unsatisfiable formula.  Repetition proves the stated
one-sided-error consequence.
\end{proof}

\begin{proposition}[Known constant rank with succinct correspondences]
\label{prop:constant-artin-rank}
Consider the following larger input model for rational Artin
motives.  Finite \'etale schemes are given by polynomial
presentations, including fixed non-split extensions
$\operatorname{Spec}\mathbb F_{q^\ell}$ and explicitly presented
fiber products with these extensions.  Pullback and pushforward along explicitly
presented finite maps, rational linear combinations,
composition, direct sums, and projector summands may be
specified by an unexpanded expression.  Validity of the
projectors is a promise.  In this model a deterministic algorithm computing
Frobenius trace modulo $p$ in time polynomial in expression
size, full rank, and $\log q$ would imply
$\mathrm{NP}=\mathrm{RP}$, even when the exact rank is given
and is the constant one for odd $p$, or the constant two for
$p=2$.
\end{proposition}
\begin{proof}
We describe all realizations explicitly, so no conjectural
motivic comparison is involved.  For a finite \'etale scheme
$Y$, its rational Artin realization is the permutation
representation on $Y(\overline{\mathbb F}_q)$; rational
idempotent correspondences define its direct summands.
Let $\pi:Y_\Phi\to\operatorname{Spec}\mathbb F_q$ be the
scheme just constructed.  The composite
\[
 \epsilon=\pi_*\pi^*:\mathbf1\longrightarrow\mathbf1
\]
is multiplication by $S_\Phi$, because pullback sends $1$ to
the constant function $1$ on all points and pushforward sums
its values.  Its expression has polynomial size without
enumerating those points.

Choose $\ell=2$ if $p$ is odd and $\ell=3$ if $p=2$, and let
$r=\ell-1$.  The permutation Artin motive of
$\operatorname{Spec}\mathbb F_{q^\ell}$ has a rational
augmentation summand $V$, obtained by subtracting its constant
projector.  It has rank $r$.  Frobenius is an $\ell$-cycle on
the full permutation representation, of trace zero; its
constant summand has trace one.  Consequently $V$ has trace
$-1$.  Its projector is
$1-\ell^{-1}\pi_\ell^*\pi_{\ell *}$, a fixed-size expression.

Here the action of the unevaluated scalar $\epsilon$ on $V$
also has an explicit allowed expression.  Write
$L=\operatorname{Spec}\mathbb F_{q^\ell}$ and let
$\pi_L:Y_\Phi\times L\to L$ be the base-changed finite map.
The polynomial presentation of this fiber product has
polynomial size, and
\[
 \epsilon_L=\pi_{L*}\pi_L^*
       =S_\Phi\operatorname{id}_{h(L)}.
\]
If $e_V$ denotes the displayed augmentation projector, then
$e_V\epsilon_L e_V$ is precisely the expression denoted
$\epsilon\operatorname{id}_V$ below.  Thus no evaluation of
$S_\Phi$ and no unspecified operation of transferring an
unevaluated scalar between objects is required.

On the promise $S_\Phi\in\{0,1\}$, both $\epsilon$ and
$1-\epsilon$ are idempotents.  Define the genuine projector
motive
\[
 M_\Phi=
 (\mathbf1^{\oplus r},(1-\epsilon)\operatorname{id})
 \ \oplus\ (V,\epsilon\operatorname{id}).
\]
If $S_\Phi=0$, it is $\mathbf1^{\oplus r}$; if $S_\Phi=1$,
it is $V$.  Thus its exact rank is always the known constant
$r$, and its trace is
\[
 \operatorname{Tr}(\operatorname{Frob}_q\mid M_\Phi)
       =r-\ell S_\Phi.
\]
The two residues are distinct modulo $p$ because $p\nmid\ell$.
All displayed expressions have polynomial size.  The same
isolation argument proves the complexity implication.
\end{proof}

The last proposition is specific to the stated succinct
correspondence syntax.  If the zero-cycle $\pi_*[Y_\Phi]$ must
instead be expanded as an integer multiple of the point, its
coefficient is already $S_\Phi$: producing that input carries
the counting problem.  Likewise a given exact rank
$S_\Phi$ trivializes the first proposition's particular
zero-dimensional example.  If the permitted objects and spans
are all split finite \'etale schemes, every resulting Artin
representation is trivial, and its trace equals its given
rank.  The non-split constant block and its base-changed spans
are therefore substantive assumptions of
Proposition~\ref{prop:constant-artin-rank}.
These distinctions prevent an
unqualified answer to the polynomial-rank question.  They
give a conditional obstruction for a broad formula model,
while leaving more restrictive, explicitly specified geometric
models to be analyzed separately.


\section{Witness degree and specialization of orbit closures}
Fix integers $1\le k\le N$.  In the affine space $W$ of degree-$N$
forms in $N^2$ variables, let $y=\det_N$ and let
$z=\ell^{N-k}\operatorname{perm}_k$ with a fixed integral embedding
of its variables.  Put $G=\mathrm{GL}_{N^2}$, the general linear group acting by linear changes of variables.  For a field $K$, let
\begin{equation}
 Y_K=\overline{G_{\overline K}\cdot y_K}\subset W_{\overline K},
 \qquad
 c_{k,N}(K)=\min\{\deg f:f\in I(Y_K),\ f(z_K)\ne0\},
\label{eq:WD1}
\end{equation}
with minimum $\infty$ if the set is empty.  This is degree in the
coefficient coordinates on $W$, not degree of the permanent itself.
The same definition applies to any fixed integral representation and
two integral points.  We write $c_0$ for characteristic zero and $c_p$
for characteristic $p$.

The algebraic closure in \eqref{eq:WD1} is essential: vanishing on the finite set
$G(\mathbb F_p)\cdot y$ would define a different problem.  However,
a minimum-degree witness may be chosen over $K$ itself.  Indeed its
existence is the rank condition below, which is unchanged by field
extension.  In the determinant representation $Y_K$ is a cone:
scalar matrices multiply a degree-$N$ form by an $N$th power, and these
powers exhaust $\overline K^*$ even when $p\mid N$.  Thus its ideal is
homogeneous, and a minimum-degree witness can also be chosen homogeneous.

\begin{lemma}[An exact bounded-degree matrix]
For every fixed $D$, there is an integer matrix $M_D$ and an integer row
$e_D$ such that over every field $K$,
\begin{equation}
 c_{k,N}(K)\le D
 \quad\Longleftrightarrow\quad
 \operatorname{rank}_K\begin{pmatrix}M_D\\e_D\end{pmatrix}
 >\operatorname{rank}_K M_D.
\label{eq:WD2}
\end{equation}
\end{lemma}
\begin{proof}
Use the substitution orbit map $g\mapsto\det_N(gx)$.  It is a
polynomial map from the space of all $N^2\times N^2$ matrices, with
integer coefficients, and its restriction to $G$ has the required
orbit.  Since $G_{\overline K}$ is dense in matrix space over every
field, a polynomial vanishes on the orbit if and only if its pullback
is the zero polynomial in the matrix entries.

Index the columns by monomials in the coefficient coordinates of degree
at most $D$.  Their pullbacks have degree at most $ND$ in the matrix
entries.  Expand them in that finite monomial basis to form $M_D$.
Let $e_D$ evaluate the original monomials at $z$.  A vector $v$ of
coefficients defines a witness precisely if $M_Dv=0$ and $e_Dv\ne0$.
Over any field this is equivalent to $e_D$ not belonging to the row
space of $M_D$, which is \eqref{eq:WD2}.  All coordinates are integral.

For a general regular $\mathrm{GL}_r$ representation, its orbit map
has entries in $\mathbb Z[g_{ij},\det(g)^{-1}]$.  For a fixed $D$,
one common power of $\det(g)$ clears every denominator.  Multiplication
by that power is injective over every field, so the identical argument
gives a finite integer coefficient matrix.
\end{proof}

\begin{samepage}
\begin{theorem}[Specialization of the minimum witness degree]
For every fixed pair $(k,N)$:
\begin{enumerate}
\item If $c_0=C<\infty$, then $c_p=C$ for all but finitely many primes.
\item If $c_0=\infty$, then for every fixed $D$, one has $c_p>D$
for all but finitely many primes.  Thus $c_p\to\infty$ in the
extended sense, allowing $c_p=\infty$.
\end{enumerate}
\end{theorem}
\end{samepage}
\begin{proof}
The rank of a fixed integer matrix equals its rational rank modulo
every prime not dividing one chosen nonzero maximal-rank minor.
All minors of larger size vanish as integers.  Applying this to both
matrices in \eqref{eq:WD2} proves that the bounded-degree existence predicate
agrees with its characteristic-zero value outside a finite, explicitly
certifiable set of primes.

If $C<\infty$, apply the assertion at $D=C$ and $D=C-1$ (or note
directly that nonzero constants cannot vanish on a nonempty orbit).
If $c_0=\infty$, apply it separately to each fixed $D$.  The exceptional
set may depend on $D$; exchanging these quantifiers would be invalid.
\end{proof}

There is a simpler certificate for the lower bound at a given $D$.
When no degree-$D$ witness exists in characteristic zero, solve
$e_D=\lambda M_D$ over $\mathbb Q$.  Clearing denominators gives
$b_De_D=\mu M_D$ with $b_D\ne0$ integral.  At any prime not dividing
$b_D$ this proves the absence of such a witness.  Conversely, clearing
denominators in one characteristic-zero witness gives an integral
polynomial whose pullback is identically zero; it remains a witness at
every prime not dividing its nonzero integer value at $z$.  These are explicit certificates for survival and bounded-degree specialization.

\begin{theorem}[Stronger stabilization for fixed orbit data]
For a fixed integral $\mathrm{GL}_r$ representation and fixed integral
points $y,z$, one has $c_p=c_0$ outside a finite set of primes, with
equality also interpreted in $\mathbb N\cup\{\infty\}$.
\end{theorem}
\begin{proof}
Only the case $c_0=\infty$ remains.  We give an explicit
scheme-image argument using elimination; in particular we do not assume
that taking an orbit closure commutes with every reduction.

Let $w_1,\ldots,w_M$ be the coordinates of the representation and write
the orbit map with coordinates in
$\mathbb Z[g_{ij},u]/(u\det(g)-1)$.  Choose polynomial representatives
$P_i(g,u)$ and form the graph ideal
\[
 J=(u\det(g)-1,\ w_i-P_i(g,u):1\le i\le M)
 \subset\mathbb Z[g_{ij},u,w_1,\ldots,w_M].
\]
Over every field this graph has coordinate ring $K[G]$, a domain;
its elimination ideal $J_K\cap K[w]$ is therefore the prime ideal
of the geometric orbit closure.

Over $\mathbb Q$, compute a finite monic Groebner basis $\mathcal G$
with an elimination order in which all $g_{ij},u$ precede the $w_i$.
Retain the finite polynomial identities expressing its members in
terms of the displayed generators of $J$, the generators in terms of
$\mathcal G$, and the reductions to zero of its finitely many
$S$-polynomials.  Invert one nonzero integer $b$ clearing every
denominator in these identities and in $\mathcal G$.  Over
$R=\mathbb Z[1/b]$ they show $J_R=(\mathcal G)$.

For every $p\nmid b$, all leading coefficients of $\mathcal G$ remain
one, its leading monomials are unchanged, and the same $S$-polynomial
reduction identities hold.  Buchberger's criterion gives a Groebner
basis $\mathcal G_p$ of $J_{\mathbb F_p}$.  The elimination theorem
therefore gives
\begin{equation}
 I(Y_{\overline{\mathbb F}_p})
   =(\mathcal H\bmod p),\qquad
 \mathcal H=\mathcal G\cap R[w],
\label{eq:WD3}
\end{equation}
after scalar extension.  The elimination order and persistence of
leading monomials ensure that no member involving an eliminated
variable suddenly becomes a $w$-only polynomial.  The same set
$\mathcal H$ generates the characteristic-zero orbit ideal.

When $c_0=\infty$, the point $z$ belongs to the characteristic-zero
orbit closure, so $h(z)=0$ in $\mathbb Q$ for every $h\in\mathcal H$.
These equalities already hold in $R$ and hence in every residue field
with $p\nmid b$.  Equation \eqref{eq:WD3} gives $z_p\in Y_p$ and $c_p=\infty$
for all such primes.  The finite-$c_0$ case was proved above.
\end{proof}

\paragraph{A second geometric proof of the closure comparison.}
Let $\mathcal Y$ be the schematic image of $G_{\mathbb Z}\to W_{\mathbb Z}$.
Its coordinate ring is the subring of $\mathbb Z[G]$ generated by the
finitely many coordinate pullbacks from $W$. It is therefore a
finite-type integral ring, torsion-free and hence flat over
$\mathbb Z$.  Its generic fiber is geometrically integral, since it
is the closure of the image of the geometrically integral group
$G_{\mathbb Q}$.  By \citet[Theorem~9.7.7(iv), pp.~79--81]{EGAIV3},
the geometrically integral fiber locus is constructible, hence contains
an open neighborhood of the generic point of $\operatorname{Spec}\mathbb Z$.
By Chevalley's theorem \citep[Theorem~1.8.4, p.~239]{EGAIV1}, the image
of $G$ in $\mathcal Y$ is constructible and contains a nonempty open
subset $U$.  This open subset meets the generic fiber.  Applying
Chevalley to $U\to\operatorname{Spec}\mathbb Z$ shows that $U_p$ is
nonempty after excluding finitely many more primes.  For each remaining
prime, $U_p$ is dense in the geometrically integral fiber $\mathcal Y_p$,
and is contained in the image of $G_p$.  After algebraic closure the
same holds: nonempty finite-type fibers over an algebraically closed
field have rational points.  Therefore the fiberwise geometric orbit
closure is exactly $\mathcal Y_p$, including its reduced scheme structure.
This proves the comparison used in \eqref{eq:WD3} independently of elimination.
No assertion of closure commutation at an excluded prime is made.

\paragraph{An exceptional prime and the model distinction.}
Let $\mathbb G_m$ act on $\mathbb A^1_{\mathbb Z}$ by scalar
multiplication, with $y=2$ and $z=1$.  In characteristic zero the orbit
of $2$ is dense in the affine line, so $c_0=\infty$.  Modulo two its
orbit is $\{0\}$, and $f(w)=w$ separates $z=1$, so $c_2=1$.
For every odd prime the orbit is again dense, so $c_p=\infty$.
The integral scheme-theoretic image is already the whole affine line,
which is flat over $\mathbb Z$; its characteristic-two fiber is not the
fiberwise orbit closure.  Flatness of the chosen closure model alone
does not remove this issue.  The finite localization in \eqref{eq:WD3} does.
Thus a single-prime witness need not lift to a characteristic-zero separation.

\paragraph{Scope as a complexity statement.}
The conclusion is for each fixed size pair $(k,N)$.  Neither the degree
$C$, the exceptional integer $b$, nor the size of the coefficient
matrices is bounded polynomially in growing $(k,N)$ by this argument.
Consequently a characteristic-zero separation premise implies finite
$c_{k,N}$ and almost-everywhere equality for each covered pair, but does
not give a uniform small witness, a uniform exceptional-prime bound,
or a circuit-size bound for a witness.  Homogeneous highest-weight
equations separating at an orbit translate may be used in characteristic
zero when assumed, but their highest-weight form and
semisimplicity need not specialize.  The theorem uses ordinary
separating polynomials and remains valid without those additional
representation-theoretic assertions.

\section{Character minors of imaginary quadratic class groups}
\label{sec:class-group-minors}
This section separates an algebraic obstruction to nonvanishing from
the analytic question of finding prime ideals in prescribed classes.
The distinction is already decisive for the ordinary ideal class group.

\begin{proposition}\label{prop:abelian-all-minors}
Let $A$ be a finite abelian group. Every square minor of its complex
character table is nonzero if and only if $|A|$ is one or a prime.
For a subset of distinct columns in a prime-order character table,
all square minors are nonzero without an analytic hypothesis.
\end{proposition}
\begin{proof}
If $|A|$ is composite, choose a proper nontrivial subgroup $H$, an
element $h\ne1$ in $H$, and a nontrivial character of $A/H$.
Pulling that character back gives $\chi\ne1$ with $\chi(h)=1$.
The rows $1,\chi$ and columns $1,h$ give the singular matrix
$\left(\begin{smallmatrix}1&1\\1&1\end{smallmatrix}\right)$.
If $|A|$ is prime, the table is the prime Fourier matrix, up to row
and column order. Chebotarev's theorem gives the conclusion; see
\citet[Lemma~1.3, author preprint pp.~3--4]{Tao2005}.
The order-one case is immediate.
\end{proof}

Here is an explicit counterexample in which both columns come from
distinct classes of unramified degree-one prime ideals.
Put $K=\mathbb Q(\sqrt{-14})$, whose integer ring is
$\mathbb Z[\sqrt{-14}]$, and write $s=\sqrt{-14}$.
The prime ideals $\mathfrak p_2=(2,s)$ and
$\mathfrak p_3=(3,s-1)$ have norms two and three.
Unique ideal factorization gives
\[
 \mathfrak p_2^2=(2),\qquad
 (2+s)=\mathfrak p_2\mathfrak p_3^2.
\]
For the second equality, the norm is $18$, the indicated primes
divide $2+s$, and its conjugate prime above three does not divide it.
The class of $\mathfrak p_2$ is nontrivial, since
$a^2+14b^2=2$ has no integer solution. It follows that
$[\mathfrak p_3]$ has order four, and that
$H=\langle[\mathfrak p_2]\rangle$ is a proper subgroup of the class group.
Choose any nontrivial character $\chi$ of the quotient by $H$.

The ideal $(3+s)$ has norm $23$ and is a principal prime.
The element $4+3s$ has norm $142=2\cdot71$, so
\[
 (4+3s)=\mathfrak p_2\mathfrak p_{71}
\]
for a prime ideal $\mathfrak p_{71}$ of norm $71$.
Thus $[\mathfrak p_{71}]=[\mathfrak p_2]\ne1$.
The two primes of norms $23$ and $71$ are unramified, have residue
degree one, and give distinct columns on which $1$ and $\chi$ agree.
Their $2\times2$ minor vanishes. It remains a minor of every table
formed using a larger cutoff $x\ge71$.
Consequently no eventual cutoff, polynomial in $\log|\operatorname{disc}K|$
or otherwise, can imply the proposed all-minors nonvanishing in general.
This counterexample also treats the interpretation in which repeated
prime ideals from the same class are first removed.

\subsection{What the analytic hypothesis would concern}
Let $\mathfrak m$ be a ray modulus and $E/K$ its ray class field.
The relevant generalized Riemann hypothesis (GRH), denoted
$\mathrm{GRH}(K,\mathfrak m)$, is that every primitive Hecke
$L_K(s,\chi)$ of finite order and
 conductor dividing $\mathfrak m$ has all nontrivial zeros on
 $\operatorname{Re}s=\tfrac12$.
The trivial character includes $\zeta_K$. Artin reciprocity identifies
these characters with those of $\Gamma=\operatorname{Gal}(E/K)$, and
\[
 \zeta_E(s)=\prod_{\chi\in\widehat\Gamma}L_K(s,\chi),
\]
with the corresponding primitive Hecke factors understood.
The Euler-factor identity can be checked directly. At a prime of $K$
with inertia subgroup $I\subseteq\Gamma$ and residue degree $f$,
only characters trivial on $I$ contribute. Their values on Frobenius
are the $f$th roots of unity, each repeated $|\Gamma|/(|I|f)$ times.
Their product is therefore
$(1-(N\mathfrak p)^{-fs})^{-|\Gamma|/(|I|f)}$, precisely the local
factor of $\zeta_E$. The nontrivial primitive Hecke factors are
holomorphic in the open critical strip, and the trivial factor's pole
is at $s=1$, so zeros in that strip cannot cancel in this product.
Thus the named hypothesis may equivalently be imposed on $\zeta_E$.
An effective Chebotarev theorem under this Dedekind-GRH hypothesis is
given by \citet[Theorem~1.1 and Corollary~1.2, corrected author version~3,
pp.~2--3]{GrenieMolteni2019}.

The scope of a Dirichlet-GRH input is explicit in the norm-induced
case. Let $\eta_K$ be the quadratic Dirichlet character of $K$, let
$\psi$ be a Dirichlet character, and let $\Sigma$ contain the rational
primes dividing their conductors or the discriminant of $K$.
Deleting the Euler factors above $\Sigma$ gives
\[
 L_K^{(\Sigma)}(s,\psi\circ N_{K/\mathbb Q})
 =L_{\mathbb Q}^{(\Sigma)}(s,\psi)
  L_{\mathbb Q}^{(\Sigma)}(s,\psi\eta_K).
\]
Indeed, a split prime $l\notin\Sigma$ gives two factors
$(1-\psi(l)l^{-s})^{-1}$ on either side. An inert prime gives
$(1-\psi(l)^2l^{-2s})^{-1}$ on the left; on the right the two
factors have opposite signs because $\eta_K(l)=-1$, giving the
same expression. The omitted finite Euler factors have zeros and
poles only on $\operatorname{Re}s=0$ and hence do not affect GRH in
the open critical strip. Dirichlet GRH for $\psi$ and $\psi\eta_K$
therefore covers this norm-induced character. No such reduction is
asserted for an arbitrary ray-class character of $K$; those characters
remain within the explicitly named additional hypothesis above.

Prime-class coverage, small generators, and nonvanishing of all
character minors are three different statements. Effective
Chebotarev concerns the first; its natural field invariants include
the ray conductor and the discriminant and degree of $E$.
Omitting an unrestricted ray conductor from a uniform bound is an
additional issue. None of these analytic estimates removes the
singular minor proved above. In prime class-group order, on the other
hand, distinct-column nonvanishing holds before any prime bound is used.

\section{Factorial divisors and weighted $p$-adic root counts}
\label{sec:div-discs}
The appropriate local quantity counts roots in discs centred at zero
and integrates over their radii. Counting all integral roots once at
each prime would give a divergent expression even for one integer root.

Let $F\in\mathbb Z[X]$ be nonzero with $F(0)\ne0$, let $c(F)>0$
be the greatest common divisor of its coefficients, and normalize $v_p(p)=1$.
Roots in $\overline{\mathbb Q}_p$ are counted with multiplicity.
For $t>0$ put
\[
 n_{F,p}(t)=\#\{\alpha:F(\alpha)=0,\ v_p(\alpha)\ge t\},
 \qquad I_p(F)=\int_0^\infty n_{F,p}(t)\,dt.
\]
The endpoints of the finitely many intervals do not affect the integral.

\begin{proposition}\label{prop:disc-product-formula}
For every such $F$,
\[
 I_p(F)=v_p(F(0))-v_p(c(F)),\qquad
 \sum_p I_p(F)\log p=\log\frac{|F(0)|}{c(F)}.
\]
If $D_N(X)=\prod_{j=1}^N(X+j)$ divides $F$ in $\mathbb Z[X]$, then
\[
 \log N!\ \le\ \sum_p I_p(F)\log p
 \ \le\ \log M(F)-\log c(F),
\]
where $M(F)=|a|\prod_\alpha\max(1,|\alpha|)$ is its complex
Mahler measure and $a$ its leading coefficient.
For the divisor itself,
\[
 I_p(D_N)=\sum_{k\ge1}\left\lfloor N/p^k\right\rfloor=v_p(N!),
 \qquad M(D_N)=N!.
\]
\end{proposition}
\begin{proof}
Write $F=a\prod_i(X-\alpha_i)$ over an algebraic closure of
$\mathbb Q_p$. The Gauss norm is multiplicative: after scaling two
polynomials to have integral coefficients with at least one unit,
their nonzero reductions multiply to a nonzero polynomial over the
residue field. Applying this to the linear factors yields
\[
 \min_jv_p([X^j]F)=v_p(a)+\sum_i\min(0,v_p(\alpha_i)).
\]
The left side is $v_p(c(F))$, whereas
$v_p(F(0))=v_p(a)+\sum_i v_p(\alpha_i)$.
Subtracting proves the first identity, since integration of the
indicator of $0<t\le v_p(\alpha_i)$ gives its positive part.
The rational product formula gives the second identity; only finitely
many primes contribute because both $F(0)$ and $c(F)$ are integers.

Write $F=D_NH$. Since $D_N$ is primitive, Gauss's lemma gives
$c(F)=c(H)$. Thus
$|F(0)|/c(F)=N!\,|H(0)|/c(H)\ge N!$.
The complex root product gives $|F(0)|\le M(F)$ directly.
Finally the roots of $D_N$ are the nonzero integers $-j$.
Their valuations are integers, and counting the multiples of $p^k$
gives the displayed floor sum. Their archimedean positive parts
give $M(D_N)=\prod_{j=1}^Nj$.
\end{proof}

\paragraph{The Newton-polygon convention.}
The root valuations used below follow from the same Gauss-norm
argument. If $F=\sum_j a_jX^j=a_d\prod_i(X-\alpha_i)$, then
for every real $t$,
\[
 \min_j\{v_p(a_j)+jt\}
 =v_p(a_d)+\sum_i\min\{t,v_p(\alpha_i)\}.
\]
For rational $t$, choose a scalar $\beta\in\overline{\mathbb Q}_p$
with $v_p(\beta)=t$ and apply Gauss multiplicativity to $F(\beta X)$.
Both sides are continuous piecewise-linear functions, so the identity
extends to real $t$. A root of valuation $v$ changes the right-hand
slope by one as $t$ crosses $v$. On the left this transition occurs
when $t$ is the negative of a slope of the lower convex hull of
the coefficient points $(j,v_p(a_j))$, with zero coefficients omitted.
Consequently the negatives of those polygon slopes are the root
valuations, and the horizontal lengths are their multiplicities.
This convention includes fractional valuations and justifies the
disc counts computed below.

If $F=\sum_{i=1}^{K}\epsilon_i f_i^M$ with $|\epsilon_i|\le1$,
and each integral $f_i$ has at most $T$ monomials of coefficient
absolute value at most $H$, then $\|F\|_1\le K(TH)^M$.
The elementary circle-mean formula for each linear
factor gives $\log M(F)=\int_0^1\log|F(e^{2\pi it})|\,dt$,
hence $M(F)\le\|F\|_1$. Therefore
\[
 \log N!\le\log K+M\log(TH)-\log c(F).
\]
This retains a coefficient-height parameter. The local identity
does not eliminate it.

\subsection{Six distinct integer roots}
Consider the sparse integer polynomials
\[
 u=360-342X+42X^3,\qquad v=119X^2-60X^3+X^4.
\]
Direct multiplication gives the complete identity
\[
 u^2-v^2=-\prod_{j=1}^6(X-j)(X^2-99X-180).
\]
In particular the polynomial is nonzero and primitive. Replacing
$X$ by $-X$ gives a rising-factorial divisor, without changing any
root valuation or the Mahler measure.
The six integer roots contribute $v_p(6!)$; the quadratic factor
has positive root valuations $(2)$ at $p=2$, $(1,1)$ at $p=3$,
and $(1)$ at $p=5$. These follow from the lower Newton polygons
of the coefficient triples $(-180,-99,1)$.
Consequently the nonzero integer-radius disc counts are
\[
\begin{array}{c|c|c|c}
 p&n_{F,p}(1)&n_{F,p}(2)&I_p(F)\\\hline
 2&4&2&6\\
 3&4&0&4\\
 5&2&0&2
\end{array}
\]
and all higher integer-radius counts vanish. Thus
\[
 \sum_p I_p(F)\log p
 =6\log2+4\log3+2\log5=\log129600=\log(6!\cdot180).
\]

\subsection{A root of multiplicity sixteen}
Now consider the sparse integer polynomials
\begin{align*}
 u={}&-57+912X^5-969X^6+969X^{16}-912X^{17}+57X^{22},\\
 v={}&272X-374X^2+646X^{10}-646X^{12}+374X^{20}-272X^{21}.
\end{align*}
Define two reciprocal polynomials by their ascending coefficient lists:
\begin{align*}
 A:\;&(57,298,846,1618,2331,2620,2331,1618,846,298,57),\\
 B:\;&(57,500,1341,2864,5036,7184,9287,10684,11366,\\
     &\hspace{20mm}10684,9287,7184,5036,2864,1341,500,57).
\end{align*}
The finite multiplication identities
\[
 u+v=(X-1)^{11}(X+1)A(X),\qquad
 u-v=(X-1)^5(X+1)B(X)
\]
give
\[
 F=u^2-v^2=(X-1)^{16}(X+1)^2A(X)B(X).
\]
The first nonzero coefficient after $X=1+Y$ is
\[
 [Y^{16}]F(1+Y)=4A(1)B(1)=4406856960
 =2^8\cdot3\cdot5\cdot11\cdot17^2\cdot19^2.
\]
It is nonzero, so the order is exactly sixteen.
This is a repeated-root example; it does not assert a divisor of
sixteen distinct consecutive roots and does not give a $16!$ bound.

To make a finite local count at the repeated root, remove that root
and put $H(Y)=F(1+Y)/Y^{16}$. It is primitive: the constant and $X^2$ coefficients of $F$ are $3249$ and $-73984$, which are coprime; integral translation preserves content, as does removal of a monomial factor.
Its positive root valuations, listed as (valuation, multiplicity), are
\[
\begin{array}{c|l|r}
 p&\text{positive valuations of roots of }H&I_p(H)\\\hline
 2&(1,4),\ (1/4,16)&8\\
 3&(1/6,6)&1\\
 5&(1/4,4)&1\\
 11&(1/10,10)&1\\
 17&(1/6,6),\ (1/12,12)&2\\
 19&(1/8,8),\ (1/14,14)&2
\end{array}
\]
For example, at two the lower hull has vertices
$(0,8),(4,4),(20,0),(28,0)$; its first two slopes give the first row.
At $3,5,11,17,19$ the negative-slope vertices are respectively
\[
\begin{gathered}
 (0,1),(6,0);\quad(0,1),(4,0);\quad(0,1),(10,0);\\
 (0,2),(6,1),(18,0);\quad(0,2),(8,1),(22,0).
\end{gathered}
\]
The remaining vertices have nonnegative slopes. These coefficient
and hull computations follow by expanding the two displayed factor
identities; they involve integer arithmetic and rational comparisons.
It follows that
\[
 \sum_p I_p(H)\log p=\log4406856960.
\]
The fractional valuations explain why a sum only over discs
$p^k\mathbb Z_p$ would lose information here. The integration is over
all positive real valuation thresholds in $\overline{\mathbb Q}_p$.

\section{Maximal affine families of odd theta characteristics}
\label{sec:odd-theta-affine}

Let \(C\) be a smooth projective geometrically connected curve of genus
\(g\geq1\) over an algebraically closed field of characteristic different
from \(2\). A theta characteristic is a line bundle whose square is
the canonical bundle. Write \(J(C)\) for the Jacobian and
\(\Theta(C)\) for the theta characteristics of \(C\),
\(V=J(C)[2]\), and \(B:V\times V\to\mathbb F_2\) for the additive Weil
pairing. Thus \(\dim_{\mathbb F_2}V=2g\), \(B\) is nondegenerate
alternating, and \(\Theta(C)\) is a \(V\)-torsor. Set
\(\epsilon(\theta)=h^0(C,\theta)\bmod2\).
The parity theorem says
\begin{equation}
 \epsilon(\theta+u+v)+\epsilon(\theta+u)
 +\epsilon(\theta+v)+\epsilon(\theta)=B(u,v).
 \label{eq:theta-polar}
\end{equation}
The number of even and odd characteristics is respectively
\(2^{g-1}(2^g+1)\) and \(2^{g-1}(2^g-1)\).
For complex curves these are Atiyah's Theorems 2 and 3, printed
pp.~48--49 \citep{Atiyah1971Spin}; Mumford's second and third results,
printed p.~182, give the algebraic statement with characteristic
different from \(2\) \citep{Mumford1971Theta}. These are the only geometric
inputs needed below.

\paragraph{Quadratic refinements and coordinates.}
For \(\theta\in\Theta(C)\) put
\[
 q_\theta(v)=\epsilon(\theta+v)+\epsilon(\theta).
\]
Equation \eqref{eq:theta-polar} makes \(q_\theta\) a quadratic refinement
of \(B\), and
\begin{equation}
 q_{\theta+u}(v)=q_\theta(v)+B(u,v).
 \label{eq:theta-translation}
\end{equation}
Nondegeneracy of \(B\) shows that \(\theta\mapsto q_\theta\) bijects the
theta characteristics with all quadratic refinements of \(B\). For a
symplectic basis \(e_i,f_i\), writing \(v=\sum_i(a_ie_i+b_if_i)\) gives
\[
 q(v)=\sum_i a_ib_i+\sum_i a_iq(e_i)+\sum_i b_iq(f_i).
\]
The elementary two-variable character sum therefore gives
\begin{equation}
 \sum_{v\in V}(-1)^{q(v)}
   =2^g(-1)^{\sum_iq(e_i)q(f_i)}.
 \label{eq:theta-gauss}
\end{equation}
This identifies the basis-independent Arf invariant. The parity counts
above imply
\(\operatorname{Arf}(q_\theta)=\epsilon(\theta)\).
In particular, for every symplectic basis there is a unique
\(\theta_0\) with \(q_{\theta_0}(e_i)=q_{\theta_0}(f_i)=0\); it is even,
and
\begin{equation}
 \epsilon\!\left(\theta_0+\sum_i(a_ie_i+b_if_i)\right)=\sum_i a_ib_i.
 \label{eq:theta-hyperbolic-coordinates}
\end{equation}
All assertions below concern this finite affine quadratic space.
They make no assertion about geometric theta characteristics of curves
in characteristic \(2\).

\begin{theorem}[Classification and exact number]
\label{thm:maximal-odd-affine}
A nonempty affine family \(A=\theta+W\subseteq\Theta(C)\) consists of
odd theta characteristics precisely when \(\epsilon(\theta)=1\) and
\(q_\theta|_W=0\). Every inclusion-maximal such family has dimension
\(g-1\), and hence cardinality \(2^{g-1}\). Moreover:
\begin{enumerate}
\item Every \((g-1)\)-dimensional \(B\)-isotropic subspace \(W\subset V\)
is the direction of exactly one maximal odd affine family.
\item These families form a single orbit under the natural symplectic
action on quadratic refinements.
\item Their number is
\begin{equation}
 M_g=(2^g-1)\prod_{j=2}^g(2^j+1).
 \label{eq:odd-family-count}
\end{equation}
Each odd theta characteristic lies in exactly
\(\prod_{j=2}^g(2^j+1)\) such families.
\end{enumerate}
Consequently, for \(g=2k\) the maximum dimension is \(2k-1\), the
maximum cardinality is \(2^{2k-1}\), and the minimum ambient codimension
is \(2k+1\).
\end{theorem}

\begin{proof}
The first criterion follows directly from the definition of \(q_\theta\).
Its polarization implies \(B|_{W\times W}=0\), so \(r=\dim W\leq g\).
Extend a basis of a \(B\)-isotropic \(r\)-space \(W\) to a symplectic
basis, with \(W=\langle e_1,\ldots,e_r\rangle\), and use
\eqref{eq:theta-hyperbolic-coordinates}. An affine coset of \(W\) has
constant parity exactly when \(b_1=\cdots=b_r=0\).
Such cosets are indexed by the remaining \(2s\) coordinates, where
\(s=g-r\), and their parity is
\(\sum_{i=r+1}^g a_ib_i\). The character sum \eqref{eq:theta-gauss}
therefore counts their odd and even constant cosets as
\begin{equation}
 O(g,r)=2^{s-1}(2^s-1),\qquad E(g,r)=2^{s-1}(2^s+1)
 \quad(s\geq1).
 \label{eq:constant-coset-count}
\end{equation}
When \(s=0\), there is one constant even coset and no odd one.
Thus no odd family has dimension \(g\), and each \(W\) of dimension
\(g-1\) has exactly one odd coset.

It remains to prove inclusion-maximality, rather than only the largest
possible dimension. Suppose \(\theta+W\) is odd and \(r<g-1\).
Here \(q_\theta\) has Arf invariant \(1\) and vanishes on \(W\).
Choose a symplectic extension \(e_i,f_i\), \(1\leq i\leq r\), of a basis
of \(W\), and replace each \(f_i\) by
\(f_i+q_\theta(f_i)e_i\). This preserves the symplectic pairings and makes
\(q_\theta(f_i)=0\). Hence these \(r\) hyperbolic planes split off an
orthogonal complement of dimension \(2s=2(g-r)\) and Arf invariant \(1\).
It has
\[
 2^{2s-1}-2^{s-1}>1
\]
zeros, because \(s\geq2\), by \eqref{eq:theta-gauss}. A nonzero singular
vector \(v\) in that complement extends \(W\) to \(W+\langle v\rangle\)
while preserving \(q_\theta=0\). Therefore every lower-dimensional odd
affine family has a one-dimensional extension. The already proved
dimension bound makes all \((g-1)\)-dimensional odd families maximal.

For completeness, count \(B\)-isotropic \(r\)-spaces by ordered bases.
After choosing an isotropic \(i\)-space, its perpendicular contains
\(2^{2g-i}\) vectors; exclude its own \(2^i\) vectors to choose the next
basis vector. Dividing by the number of ordered bases of an \(r\)-space
gives
\begin{equation}
 I(g,r)=
 \prod_{i=0}^{r-1}
 \frac{2^{2g-i}-2^i}{2^r-2^i}
 =
 \frac{\prod_{j=g-r+1}^g(2^{2j}-1)}
      {\prod_{j=1}^r(2^j-1)}.
 \label{eq:isotropic-space-count}
\end{equation}
At \(r=g-1\) this simplifies to \eqref{eq:odd-family-count}.
The symplectic group is transitive on isotropic subspaces of fixed
dimension, by extension to symplectic bases. Its action
\(q\mapsto q\circ S^{-1}\) preserves the Arf invariant and transports
the unique odd coset attached to a direction; this proves the orbit
claim. This action is also transitive on Arf-invariant-one refinements:
successively split off hyperbolic planes using a nonzero singular vector
until only a two-dimensional anisotropic plane remains, which has a
unique nonzero-value quadratic form. Double-counting incident odd
points and maximal families gives
\[
 \frac{M_g\,2^{g-1}}{2^{g-1}(2^g-1)}
 =\prod_{j=2}^g(2^j+1)
\]
families through each odd point.
\end{proof}

\paragraph{Classical terminology and scope.}
An isotropic group of characteristics is a G\"opel group; its affine
cosets are G\"opel systems. The constant-coset counts
\eqref{eq:constant-coset-count}, including uniqueness of the odd coset
at rank \(g-1\), are explicitly recorded in
Bernatska--Kopeliovich, Section~2.6, p.~7
\citep{BernatskaKopeliovich2023}, which refers to Baker's Articles
297 and 299--300. Thus the classification is classical theta
combinatorics, and no novelty is claimed.

One can state exactly the relevant meaning of \emph{syzygetic}.
A triple \(\theta_1,\theta_2,\theta_3\) is syzygetic when
\[
 \epsilon(\theta_1)+\epsilon(\theta_2)+\epsilon(\theta_3)
 +\epsilon(\theta_1+\theta_2+\theta_3)=0;
\]
the sum of three points in an \(\mathbb F_2\)-affine torsor is intrinsic.
Every triple in an odd affine family is syzygetic. Conversely, if a
nonempty set \(S\) of odd characteristics has all its triples
syzygetic, fix \(\theta\in S\). The differences \(u,v\) from \(\theta\)
satisfy \(q_\theta(u)=q_\theta(v)=q_\theta(u+v)=0\), so \(B(u,v)=0\).
Their linear span is therefore \(q_\theta\)-singular, and the affine
span of \(S\) is entirely odd. Thus the maximal such syzygetic sets are
exactly the maximal affine families in the theorem. This statement
does not require theta characteristics to have exactly one section,
nor does it assert a configuration theorem for bitangent hyperplanes.

\begin{corollary}[Exact affine covering number]
\label{cor:odd-theta-cover}
Covering all odd theta characteristics by affine families contained
in the odd locus requires and suffices to use \(2^g-1\) families.
More generally, if such a cover contains at least a fraction \(\delta\)
of all odd characteristics, it has at least
\(\lceil\delta(2^g-1)\rceil\) members.
\end{corollary}
\begin{proof}
Each family has at most \(2^{g-1}\) points, whereas the odd locus has
\(2^{g-1}(2^g-1)\) points. This proves the lower bounds by the union
bound. For the upper bound, use \(q(a,b)=a\cdot b\) and, for each
nonzero \(b\in\mathbb F_2^g\), take
\[
 A_b=\{(a,b):a\cdot b=1\}.
\]
The \(2^g-1\) sets \(A_b\) are disjoint affine \((g-1)\)-flats and
partition the odd locus.
\end{proof}

\paragraph{Application to inner-product parity fibres.}
For \(k\geq1\) four-bit inner-product blocks put
\[
 P(y)=\sum_{i=1}^k(y_{i,0}y_{i,1}+y_{i,2}y_{i,3}).
\]
This is the hyperbolic quadratic form of genus \(g=2k\). If a
nonempty affine fibre \(F\subseteq\mathbb F_2^{4k}\) satisfies
\(P|_F=1\), then
\begin{equation}
 \dim F\leq2k-1,\qquad \operatorname{codim}F\geq2k+1.
 \label{eq:odd-ip-rank}
\end{equation}
Both inequalities are sharp. The even fibre has maximum affine
dimension \(2k\), so the bound there is only codimension \(2k\), also
sharp. In resolution over parity equations, denoted Res\((\oplus)\),
a line is a disjunction of affine equations over \(\mathbb F_2\).
Its rank is the rank of their coefficient matrix; its falsifying set
is the affine solution set obtained by negating every equation.
If \(F\) is obtained by fixing the outside coordinates of the
falsifying affine flat of a line of rank \(r\), its codimension is at
most \(r\); thus \(P|_F=1\) forces \(r\geq2k+1\).

This sharpens the isotropic rank estimate by one for an odd
wrong-parity fibre. A boundary condition \(P=c\) has wrong-parity fibre
\(P=1\) only when \(c=0\); when \(c=1\), the wrong-parity fibre is even
and the improvement by one does not apply. Moreover,
Corollary~\ref{cor:odd-theta-cover} is a cardinality statement with
respect to the uniform distribution on the full odd fibre. It supplies
neither a common boundary set for arbitrary proof lines nor a lower
bound on the mass of their union under a
distribution used in a proof-complexity argument. For a refutation $\pi$, a distribution $\mu$, falsifying sets $A_L$ of its lines, a rank threshold $\rho$, and a positive constant $\delta$, the required assertion would have the form
\begin{equation}
 \sum_{\substack{L\in\pi\\\operatorname{rank}L\geq\rho}}
       \mu(A_L)\geq\delta.
 \label{eq:theta-m3-not-proved}
\end{equation}
Establishing this covering inequality requires an additional argument. No
superpolynomial Res\((\oplus)\) lower bound follows without it.

\paragraph{Finite examples.}
The accompanying integer-bit program enumerates all symplectic
isotropic directions through genus \(4\), all odd cosets in each
direction, and every one-step extension of a lower-dimensional odd
flat. Through genus \(3\) it additionally enumerates all affine
directions, without imposing isotropy, and obtains the same odd
families. The resulting maximal-family counts are
\[
\begin{array}{c|rrrr}
g&1&2&3&4\\ \hline
\#\Theta_{\mathrm{odd}}&1&6&28&120\\
\text{maximum affine dimension}&0&1&2&3\\
\#\text{maximal families}&1&15&315&11475 .
\end{array}
\]
These finite checks illustrate the preceding proof; the proof applies in every genus.



\section{What remains open}
The second height stratum of the counting hypersurfaces is characterized
and modular-counting complete at every fixed prime. The full height
has a terminating coefficient rule and a polynomial-space upper bound.
With the prime included in the input, it is also $\#\mathrm P$-hard
under nonadaptive polynomial-time oracle reductions.
A sharp complexity classification of the unbounded exact-height
function remains open. Exact-height decision with a growing unary
height parameter is modular-counting complete at a fixed prime; the
proof gives no such classification for a binary height parameter or
for infinite-height decision. The replica formula ceases to be the entire
coefficient once $p^j-1\ge Q(Q+1)$; understanding those additional
selections may lead to more economical tests or to further fixed-height
gadgets. No determination of all higher strata from $S$ is asserted.

The bounded-rank trace obstruction concerns the expressly stated
succinct correspondence syntax. It does not decide polynomial trace
complexity for every geometric presentation with an expanded basis
or expanded correspondences. Witness-degree specialization supplies no
polynomial bound as determinant and permanent sizes grow. The theta
classification supplies a local dimension and cardinality bound,
but the distributional covering step for arbitrary refutations is
still required.

\section*{Method and records}
The mathematics and text were produced by GPT-6 Astra (ChatGPT Pro mode). Separate re-proofs by the same model, full mathematical readings, and exact integer and finite-field computations are distinguished in the check records at \url{\recordurl}. No human has checked this paper. The programs and certificates in the workbench at \url{\workbenchurl} include \emph{Replica coefficient verification}, \emph{Projective point-count precision}, \emph{Third-height inclusion--exclusion}, \emph{Witness-degree matrix checks}, \emph{Artin-rank checks}, \emph{Valuation disc calculations}, and \emph{Odd affine families}. This text is dedicated to the public domain under Creative Commons CC0~1.0.
\appendix
\section{An alternative polynomial-space proof from finite traces}
\label{subsec:height-polynomial-space}
The main resource theorem follows from the quasi-$F$-splitting matrix. This appendix gives a second proof from Lauder and Wan's finite traces. The finite coefficient rule settles termination, but its largest exponent is too large for a direct space bound. We instead compute the relevant Frobenius coefficients through point counts. Lauder and Wan's finite trace formula permits these counts to be represented by arithmetic circuits with polynomial depth. Uniform integer division then makes their bits accessible without storing the exponentially long integers themselves.

\begin{theorem}[Exact height in polynomial space]
\label{thm:height-polynomial-space}
Let a binary-encoded prime $p$, a division-free arithmetic circuit over $\mathbb F_p$, and its $d\ge3$ explicitly listed variables be given. Promise that the circuit represents a homogeneous degree-$d$ polynomial defining a smooth hypersurface $X\subset\mathbb P^{d-1}$. The function returning its top Artin--Mazur height, as a binary integer or the symbol $\infty$, belongs to $\mathrm{FPSPACE}$.
\end{theorem}

Here $\mathrm{FPSPACE}$ means deterministic computation using a number of work bits polynomial in input length. We prove the theorem through three lemmas, including the numerical precision and the uniformity needed for this model.

\begin{lemma}[The first nondivisible Frobenius coefficient]
\label{lem:height-first-frobenius-coefficient}
Let $d\ge4$, $n=d-2$, and let
\[
 P(T)=\det(1-T\operatorname{Frob}_p\mid H^n_{\mathrm{prim}}(X))
      =\sum_{j\ge0}c_jT^j\in1+T\mathbb Z[T].
\]
Write $c_j=0$ beyond its degree. With $R=B_d$ as in \eqref{eq:HC1}, the height is the first $j\le R$ such that $p^j\nmid c_j$. If no such index exists, the height is infinite.
\end{lemma}
\begin{proof}
Section~\ref{sec:point-count-height-two} identifies the reciprocal roots of $P$ with the primitive crystalline Frobenius eigenvalues and proves its integrality. If the finite height is $h$, exactly $h$ roots have valuation $(h-1)/h$ and every other root has valuation at least one. For $j<h$, every term of the $j$th elementary symmetric coefficient has valuation at least
\[
 j(1-1/h)>j-1.
\]
The integral coefficient therefore has valuation at least $j$. At $j=h$, the product of the $h$ small-slope roots is the unique term of valuation $h-1$; every other product includes a root of valuation at least one and has larger valuation. Thus $v_p(c_h)=h-1$. This also covers $h=1$. If height is infinite, all roots have valuation at least one and $p^j\mid c_j$ for every $j$. The finite-height bound makes the cutoff $R$ sufficient.
\end{proof}

\begin{lemma}[Arithmetic circuits and space]
\label{lem:implicit-arithmetic-space}
Suppose an integer-valued construction on inputs of length $L$ has Boolean circuits of size at most $2^{\operatorname{poly}(L)}$ and depth $\operatorname{poly}(L)$. Suppose also that a gate and its predecessors can be described using polynomially many bits and computed in polynomial space. Then any output bit is computable in polynomial space.

The required circuit bounds hold for addition, subtraction, multiplication, comparison, integer quotient and remainder on operands with $2^{\operatorname{poly}(L)}$ bits. They also hold for balanced sums and products of exponentially many operands, and for repeated squaring with an exponent of polynomial binary length, provided all intermediate bit lengths remain exponentially bounded.
\end{lemma}
\begin{proof}
Evaluate a gate recursively, recomputing predecessors when needed. For bounded fan-in, each stack level stores its gate identifier, which predecessor is being evaluated, and a bounded number of returned bits. Polynomial depth gives a polynomial stack. The space used to generate a predecessor is reused at each call. An unbounded threshold gate can instead be evaluated by scanning its predecessors and maintaining a counter of polynomial binary length.

Carry-lookahead addition and balanced summation of partial products give uniform polynomial-size, polylogarithmic-depth circuits for ordinary integer arithmetic except division. Hesse, Allender and Barrington give division in DLOGTIME-uniform $\mathrm{TC}^0$, meaning constant-depth, polynomial-size threshold circuits whose local wiring is computable in deterministic logarithmic time \citep[Corollary~6.3, author manuscript p.~16; uniformity conventions pp.~4--6]{HAB2002}. A threshold gate can be replaced by balanced binary addition and comparison, so bounded-fan-in circuits of polylogarithmic depth suffice. Their corrigendum retracts two consequences in Corollary~6.6 and explicitly preserves the division theorem \citep[pp.~1--2]{HABCorr2014}. No retracted consequence is used here. Quotient gives remainder by multiplication and subtraction; signs are handled separately.

The arithmetic circuit descriptions are locally computable in logarithmic space in the operand length. This is polynomial space when that length is exponentially bounded. Balanced products add logarithmically many arithmetic stages, and repeated squaring adds logarithmically many stages in the exponent. These are polynomially many stages under the stated hypotheses.
\end{proof}

\begin{lemma}[Point-count bits at bounded extension indices]
\label{lem:point-count-bits-space}
Under the hypotheses of Theorem~\ref{thm:height-polynomial-space}, every bit of
\[
 N_k=\#X(\mathbb F_{p^k}),\qquad 1\le k\le B_d,
\]
has an exponential-size, polynomial-depth Boolean circuit with polynomial-space access to its gates. The bounds are polynomial in the original input length, including the binary lengths of $p$ and $k$.
\end{lemma}
\begin{proof}
Put $R=B_d$. Since $R\le\max(36,4^d)$, its binary length is $O(d)$. Partition projective space by its nonempty coordinate supports. On each support normalize its first nonzero coordinate to one. The resulting $2^d-1$ disjoint strata are tori of dimensions $m\le d-1$, with one restricted equation $f$ of degree at most $d$.

Each coefficient of $f$ has the required circuit description. Substitute the prescribed zeros and one into the original circuit and evaluate in the polynomial ring over $\mathbb F_p$ modulo the ideal of monomials of total degree greater than $d$. There are at most $(d+1)^d$ coefficient positions. A product coefficient is a balanced convolution, and there are only as many successive stages as original circuit gates. Truncation is exact because all exponents are nonnegative. The promise identifies the output with the desired equation of degree at most $d$. Balanced tests of its coefficient array distinguish zero and constant equations and determine its degree $\delta$. Equivalently, construct the $d$ possible positive-degree branches and select the appropriate one. A zero equation contributes $(p^k-1)^m$, a nonzero constant contributes zero, and a zero-dimensional torus is immediate.

\paragraph{The finite trace formula.}
For a nonconstant equation of degree $\delta$ in $m\ge1$ torus variables, use Lauder and Wan's construction with unramified degree $a=1$. Thus its coefficient Frobenius automorphism is the identity. Set
\[
 N=(m+1)k,\qquad
 R_N=(\mathbb Z/p^N\mathbb Z)[\pi]/(\pi^{p-1}+p),\qquad
 t=\left\lceil\frac{p^2N}{(p-1)^2}\right\rceil-1<4N.
\]
Write $X_0f=\sum_j\bar a_jX^j$, let $a_j$ be the Teichm\"uller lift of $\bar a_j$ modulo $p^N$, and put
\[
 \theta(z)=\exp(\pi z-\pi z^p),\qquad
 \mathcal F(X)=\prod_j\theta(a_jX^j).
\]
Here every $j_0=1$. Zero coefficients may be included, with their factors equal to one.

For the degree-$\delta$ simplex, the matrix $A_N$ of Lauder--Wan's Theorem~28, pp.~600--601 \citep{LauderWan2008}, has indices
\[
 \mathcal U=\{u\in\mathbb Z_{\ge0}^{m+1}:u_0\le t,
                         \ u_1+\cdots+u_m\le\delta u_0\}
\]
and entries
\[
 (A_N)_{u,v}=[X^{pu-v}]\mathcal F\pmod{p^N},
\]
with negative exponents giving zero. Let $N_k^*$ be the toric count. The theorem gives
\[
 (p^k-1)^{m+1}\operatorname{Tr}(A_N^k)
 \equiv p^kN_k^*-(p^k-1)^m\pmod{p^N}.
\]
Consequently Algorithm~29, pp.~602--603, reconstructs it exactly as
\begin{equation}\label{eq:space-toric-count}
 N_k^*=\frac1{p^k}
 \left[(p^k-1)^{m+1}\operatorname{Tr}(A_N^k)
                         +(p^k-1)^m\right]_{[0,p^N)}.
\end{equation}
The bracket is the least nonnegative integer residue. The trace formula ensures that the ring expression is an integer residue: read its coefficient of $\pi^0$ in the basis $1,\pi,\ldots,\pi^{p-2}$. All other components vanish. Moreover
\[
 0\le p^kN_k^*\le p^k(p^k-1)^m<p^{(m+1)k}=p^N,
\]
so the residue is the integer itself and the final division is exact. This specifies the precision; no reconstruction ambiguity remains.

For regular addressing, pad the index set to
\[
 0\le u_0\le t,\qquad 0\le u_i\le dt\quad(1\le i\le m),
\]
putting zero in every row and column outside $\mathcal U$. Zero padding preserves traces of positive powers. Its order is
\[
 W=(t+1)(dt+1)^m.
\]
All entries needed in the matrix have first exponent at most $K=pt$ and every other exponent at most $dK$. It therefore suffices to compute the coefficients of $\mathcal F$ in this rectangle. We do not need its entire power series.

\paragraph{Lifting at polynomial depth.}
The full representation of an element of $R_N$ can be exponentially long, but it has only $p-1$ components, each with $\lceil N\log_2p\rceil$ bits. Both component and bit addresses have polynomial length. Multiplication is a convolution in these components followed by $\pi^{p-1}=-p$ and reduction modulo $p^N$. The convolution sums are balanced.

For a coefficient $c\in\mathbb F_p$, Newton-lift the equation
\[
 \phi(z)=z^p-z,
 \qquad \phi'(z)=pz^{p-1}-1.
\]
Start with $g=c$ and $s=-1\pmod p$. If $\phi(g)=0\pmod{p^b}$ and $s\phi'(g)=1\pmod{p^b}$, define at doubled precision
\[
 g'=g-\phi(g)s,\qquad s'=2s-\phi'(g')s^2.
\]
Taylor expansion gives $\phi(g')=0\pmod{p^{2b}}$, and
$1-s'\phi'(g')=(1-s\phi'(g'))^2$ gives the inverse-derivative invariant. Reduce at precision $p^{\min(2b,N)}$. After $O(\log N)$ rounds, uniqueness of the lift with the specified residue gives the Teichm\"uller coefficient. This includes $c=0$ and $p=2$. Each evaluation of $z^p$ requires only $O(\log p)$ multiplication stages. This is the quadratic-convergence method of Lauder--Wan Lemma~5, pp.~586--587, with its depth made explicit.

\paragraph{Splitting coefficients without a long recurrence.}
For $r\le K$ put
\[
 \alpha_r=v_p(r!)=\sum_{i\ge1}\lfloor r/p^i\rfloor,
 \quad u_r=r!/p^{\alpha_r},
 \quad e_r=r-(p-1)\alpha_r\ge0.
\]
The last expression is the base-$p$ digit sum of $r$. Thus
\begin{equation}\label{eq:space-integral-exponential}
 E_r:=\frac{\pi^r}{r!}
       =(-1)^{\alpha_r}\pi^{e_r}u_r^{-1}\in\mathbb Z_p[\pi].
\end{equation}
Compute $r!$ by a balanced product, the valuation by its logarithmically many floor terms, and $u_r$ by exact integer division. The factorial has $O(K\log(K+1))$ bits. Since $u_r$ is a unit modulo $p$, its inverse modulo $p$ is obtained by powering to $p-2$. Precision doubling
\[
 v\longmapsto v(2-u_rv)
\]
then gives its inverse modulo $p^N$ in $O(\log N)$ rounds. Formula~\eqref{eq:space-integral-exponential} divides by no nonunit in $R_N$; powers of $\pi$ reduce using $\pi^{p-1}=-p$.

Multiplication of the two formal exponentials gives the finite sum
\[
 \theta_r=[z^r]\theta(z)
   =\sum_{j=0}^{\lfloor r/p\rfloor}(-1)^jE_{r-pj}E_j.
\]
Every summand is integral by \eqref{eq:space-integral-exponential}; the sum is balanced. This is the coefficient construction underlying Lauder--Wan Lemma~30, p.~604. It avoids a recurrence with $K$ successive stages. In characteristic two the ring has one basis component and $\pi=-2$, so the same formulas apply.

For a factor $\theta(a_jX^j)$, each $a_j^r$ uses $O(\log(r+1))$ squaring stages. Multiply the at most $(d+1)^d$ factors in a balanced tree, with coordinate truncation at $K,dK,\ldots,dK$. Because exponents are nonnegative, a discarded term cannot affect a coefficient inside that rectangle. Its number of positions is
\[
 V=(K+1)(dK+1)^m.
\]
Each polynomial multiplication uses a balanced sum of at most $V$ ring products per coefficient. The matrix entries follow by index selection. Compute $A_N^k$ by repeated squaring, with balanced dot products, and take its trace by balanced summation. Formula~\eqref{eq:space-toric-count} then uses the arithmetic operations of Lemma~\ref{lem:implicit-arithmetic-space}.

\paragraph{Uniform size and depth.}
If $s$ is the input circuit size, set
\[
 A=2+d+s+\lceil\log_2p\rceil+\lceil\log_2(R+1)\rceil=O(L).
\]
The parameters $N,t,K$ and all integer bit widths above are $2^{O(A)}$; the array dimensions $V,W$ are $2^{O(A^2)}$. There are only polynomially many successive stages: the original $s$ circuit stages; $O(\log N)$ lifting rounds; logarithmically many powering, factor-product and summation stages; and $O(\log k)$ matrix-power stages. Each stage has polynomial Boolean depth in the logarithms of these dimensions and widths. The total size is $2^{\operatorname{poly}(L)}$ and depth is $\operatorname{poly}(L)$.

Gate descriptions contain the support mask, degree branch, coefficient exponent vector or matrix indices, ramified basis index, arithmetic bit address, and balanced-tree or iteration index. There are polynomially many coordinates, each of polynomial binary length. Their predecessor descriptions use integer arithmetic on these addresses and the locally uniform circuits of Lemma~\ref{lem:implicit-arithmetic-space}. Zero and constant tests and the degree selection are Boolean gates, rather than unexplained preprocessing. Thus the required gate access is in polynomial space. Summing the toric counts over all supports adds exponentially many independent copies and only polynomial depth, proving the lemma.
\end{proof}

\begin{proof}[Proof of Theorem~\ref{thm:height-polynomial-space}]
First assume $d\ge4$, put $n=d-2$, and set $R=B_d$. The primitive power sums are
\[
 \tau_k=(-1)^n\left(N_k-\sum_{i=0}^{n}p^{ik}\right),
 \qquad 1\le k\le R.
\]
Lemma~\ref{lem:point-count-bits-space} gives their bit circuits. The ambient contribution, including its slope-zero term, has been removed. Their bit lengths are exponential at most, since
\[
 |\tau_k|\le 2d\,p^{(d-1)k}.
\]
We next obtain the coefficients of
\[
 P(T)\equiv\exp\left(-\sum_{k=1}^{R}\tau_kT^k/k\right)
                  \pmod{T^{R+1}}
\]
without a sequential recurrence of length $R$.

Put $B=R!$ and define integer polynomials
\[
 Q_k(T)=\sum_{r=0}^{\lfloor R/k\rfloor}
        \frac{B(-\tau_k)^r}{k^r r!}T^{kr},
 \qquad 1\le k\le R.
\]
Their coefficients are integral: $(kr)!/(k^rr!)$ counts permutations of $kr$ letters consisting of $r$ cycles of length $k$, so $k^rr!$ divides $(kr)!$ and hence divides $R!$. Compute
\[
 Q(T)=\prod_{k=1}^{R}Q_k(T)\pmod{T^{R+1}}
\]
in a balanced tree. The formal exponential identity gives
\begin{equation}\label{eq:space-frobenius-reconstruction}
 c_j=[T^j]Q/B^R\qquad(0\le j\le R).
\end{equation}
This division is exact because $P$ has integer coefficients. Factorials and powers in each $Q_k$ use balanced products and repeated squaring.

For the bit bounds, $\log_2B\le R\log_2(R+1)$, and each coefficient of one $Q_k$ has
\[
 O\bigl(R^2d\log p+R\log(R+1)\bigr)
\]
bits. A coefficient of a product of at most $R$ such factors is a sum of at most $(R+1)^R$ products. Consequently all intermediate coefficients, and the denominator $B^R$, have at most
\[
 O\bigl(R^3d\log p+R^2\log(R+1)\bigr)=2^{\operatorname{poly}(L)}
\]
bits. There are only $O(\log R)$ polynomial-product stages. Thus \eqref{eq:space-frobenius-reconstruction}, integer division and reduction modulo $p^j$ give the same succinct circuit bounds for every desired remainder.

Enumerate $j=1,\ldots,R$ with a polynomial-length counter. Test the remainder of $c_j$ modulo $p^j$ one bit at a time using Lemma~\ref{lem:implicit-arithmetic-space}. The modulus itself has an exponential number of bits, but its exponent $j$ has polynomial binary length and its bit circuit uses polynomially many squaring stages. Return the first nonzero remainder's index, or $\infty$ if all vanish. Lemma~\ref{lem:height-first-frobenius-coefficient} proves correctness. The counters, recursive circuit-evaluation stack and gate-generation workspace are all polynomially bounded.

For $d=3$, use the coefficient-height test directly with $B_3=2$. Form the homogeneous integral component-circuit lift as above, compute $a_p\bmod p$ and, if it vanishes, $a_{p^2}\bmod p^2$. These distinguish heights one, two and infinity. For both powers, evaluate in the polynomial ring truncated at exponent $p^2$ in each of the three variables. Its array has at most $p^6$ positions; the scalar modulus is $p^2$ and the powers require $O(\log p)$ squaring stages. Balanced convolutions and the polynomial number of original circuit stages again give exponential size and polynomial depth. Hence the same polynomial-space bit evaluator handles this case without a separate crystalline comparison for curves.

The output is at most $R$ or the single exceptional symbol, so it has polynomial length. This completes the uniform deterministic space bound.
\end{proof}

\begin{corollary}[Deterministic exponential time]
\label{cor:height-exponential-time}
The exact-height function, including infinite height and a binary-encoded input prime, belongs to $\mathrm{FEXPTIME}$.
\end{corollary}
\begin{proof}
A halting deterministic polynomial-space computation with polynomial output length has at most exponentially many distinct configurations, after buffering that output in polynomial workspace. It cannot repeat a configuration. The preceding algorithm therefore runs in $2^{\operatorname{poly}(L)}$ time. Lauder and Wan's zeta algorithm also gives this time bound directly \citep[Theorems~1 and~37, pp.~580 and~610]{LauderWan2008}; the theorem above additionally controls space.
\end{proof}


\setlength{\bibsep}{3pt}
\begin{thebibliography}{99}
\bibitem[Artin and Mazur(1977)]{AM} Artin, M., \& Mazur, B. (1977). Formal groups arising from algebraic varieties. \emph{Annales scientifiques de l'\'Ecole Normale Sup\'erieure, 10}(1), 87--131. \url{https://doi.org/10.24033/asens.1322}.
\bibitem[Atiyah(1971)]{Atiyah1971Spin} Atiyah, M. F. (1971). Riemann surfaces and spin structures. \emph{Annales scientifiques de l'\'Ecole Normale Sup\'erieure, 4}(1), 47--62. \url{https://doi.org/10.24033/asens.1205}.
\bibitem[Batubara et~al.(2026)]{BGHM2026} Batubara, R., Garzella, J. J., Huang, Y., \& Mellberg, M. (2026). Newton strata realization for hypersurfaces via explicit $p$-adic cohomology. \emph{Research in Number Theory, 12}, Article 110. \url{https://doi.org/10.1007/s40993-026-00802-6}.
\bibitem[Batubara et~al.(2025)]{BGP2025} Batubara, R., Garzella, J. J., \& Pan, A. (2025). \emph{K3 surfaces of any Artin--Mazur height over $\mathbb F_5$ and $\mathbb F_7$ via quasi-$F$-split singularities and GPU acceleration} (Version 3). arXiv. \url{https://arxiv.org/abs/2502.12428v3}.
\bibitem[Beigel and Gill(1992)]{BeigelGill1992} Beigel, R., \& Gill, J. (1992). Counting classes: Thresholds, parity, mods, and fewness. \emph{Theoretical Computer Science, 103}(1), 3--23. \url{https://doi.org/10.1016/0304-3975(92)90084-S}. Author manuscript, Theorem~30, p.~15.
\bibitem[Bernatska and Kopeliovich(2023)]{BernatskaKopeliovich2023} Bernatska, J., \& Kopeliovich, Y. (2023). \emph{Igusa's map on modular forms vanishing on the hyperelliptic locus}. arXiv:2306.14889v1, Section~2.6, p.~7. \url{https://arxiv.org/abs/2306.14889v1}.
\bibitem[Consequences, corrected transfers, and remaining gates(2026)]{ParityCompanion} \emph{Consequences, corrected transfers, and remaining gates}. (2026). Unpublished manuscript, 22 November 2026; cited by subtitle. Complexity chapter, subsections ``Parity hardness persists when the coherent Frobenius matrix has only one entry'' and ``The parity gate survives smooth hypersurfaces and genuine finite Weil objects''. Records at \url{\recordurl}.
\bibitem[Counting satisfying assignments(2026)]{Counting} \emph{Counting satisfying assignments on smooth projective varieties, after Katz}. (2026). Unpublished companion manuscript, 24 November, Theorems~1.1--1.2 and Sections~2--3. Records at \url{\recordurl}.
\bibitem[Deligne(1974)]{Deligne1974} Deligne, P. (1974). La conjecture de Weil: I. \emph{Publications Math\'ematiques de l'IH\'ES, 43}, 273--307. \url{https://doi.org/10.1007/BF02684373}.
\bibitem[Greni\'e and Molteni(2019)]{GrenieMolteni2019}
Greni\'e, L., \& Molteni, G. (2019).
An explicit Chebotarev density theorem under GRH.
\emph{Journal of Number Theory, 200}, 441--485.
\url{https://doi.org/10.1016/j.jnt.2018.12.005}.
Corrected author version: arXiv:1709.07609v3,
Theorem~1.1 and Corollary~1.2, pp.~2--3.
\url{https://arxiv.org/abs/1709.07609v3}.
\bibitem[Grothendieck and Dieudonn\'e(1964)]{EGAIV1} Grothendieck, A., \& Dieudonn\'e, J. (1964). \emph{\'El\'ements de g\'eom\'etrie alg\'ebrique: IV, premi\`ere partie}. Publications Math\'ematiques de l'IH\'ES, 20. Theorem~1.8.4, p.~239.
\bibitem[Grothendieck and Dieudonn\'e(1966)]{EGAIV3} Grothendieck, A., \& Dieudonn\'e, J. (1966). \emph{\'El\'ements de g\'eom\'etrie alg\'ebrique: IV, troisi\`eme partie}. Publications Math\'ematiques de l'IH\'ES, 28. Theorem~9.7.7, pp.~79--81.
\bibitem[Hesse et~al.(2002)]{HAB2002} Hesse, W., Allender, E., \& Barrington, D. A. M. (2002). Uniform constant-depth threshold circuits for division and iterated multiplication. \emph{Journal of Computer and System Sciences, 65}(4), 695--716. \url{https://doi.org/10.1016/S0022-0000(02)00025-9}. Corollary~6.3, author manuscript p.~16: \url{https://people.cs.rutgers.edu/~allender/papers/division.pdf}.
\bibitem[Hesse et~al.(2014)]{HABCorr2014} Hesse, W., Allender, E., \& Barrington, D. A. M. (2014). Corrigendum to ``Uniform constant-depth threshold circuits for division and iterated multiplication.'' \emph{Journal of Computer and System Sciences, 80}(2), 496--497. \url{https://doi.org/10.1016/j.jcss.2013.09.002}. Author manuscript pp.~1--2: \url{https://people.cs.rutgers.edu/~allender/papers/hab.corrigendum.pdf}.
\bibitem[Katz(1972)]{Katz1972} Katz, N. M. (1972). Algebraic solutions of differential equations ($p$-curvature and the Hodge filtration). \emph{Inventiones Mathematicae, 18}, 1--118. Algorithm~2.3.7.14 and Special Case~2.3.7.17, pp.~35--36.
\bibitem[Katz and Messing(1974)]{KatzMessing1974} Katz, N. M., \& Messing, W. (1974). Some consequences of the Riemann hypothesis for varieties over finite fields. \emph{Inventiones Mathematicae, 23}, 73--77. Theorem~1, p.~74. \url{https://doi.org/10.1007/BF01405203}.
\bibitem[Kawakami et~al.(2025)]{KTY2025} Kawakami, T., Takamatsu, T., \& Yoshikawa, S. (2025). \emph{Fedder type criteria for quasi-$F$-splitting I} (Version 3; original preprint 2022). arXiv. \url{https://arxiv.org/abs/2204.10076v3}. Theorem~C already appears in version~2, p.~5: \url{https://arxiv.org/abs/2204.10076v2}.
\bibitem[Lauder and Wan(2008)]{LauderWan2008} Lauder, A. G. B., \& Wan, D. (2008). Counting points on varieties over finite fields of small characteristic. In J. P. Buhler \& P. Stevenhagen (Eds.), \emph{Algorithmic number theory: Lattices, number fields, curves and cryptography} (pp.~579--612). Cambridge University Press. \url{https://library.slmath.org/books/Book44/files/18lauder.pdf}.
\bibitem[Milne(2013)]{MilneLEC2013} Milne, J. S. (2013). \emph{Lectures on \'etale cohomology} (Version~2.21). Theorems~15.1 and~24.1, pp.~105 and~144. \url{https://www.jmilne.org/math/CourseNotes/LEC.pdf}.
\bibitem[Mumford(1971)]{Mumford1971Theta} Mumford, D. (1971). Theta characteristics of an algebraic curve. \emph{Annales scientifiques de l'\'Ecole Normale Sup\'erieure, 4}(2), 181--192. \url{https://doi.org/10.24033/asens.1209}.
\bibitem[Stienstra(1987)]{Stienstra} Stienstra, J. (1987). Formal group laws arising from algebraic varieties. \emph{American Journal of Mathematics, 109}(5), 907--925. \url{https://doi.org/10.2307/2374494}.
\bibitem[Tao(2005)]{Tao2005} Tao, T. (2005). An uncertainty principle for cyclic groups of prime order. \emph{Mathematical Research Letters, 12}(1), 121--127. Lemma~1.3. \url{https://doi.org/10.4310/MRL.2005.v12.n1.a11}.
\bibitem[Valiant and Vazirani(1986)]{ValiantVazirani1986} Valiant, L. G., \& Vazirani, V. V. (1986). NP is as easy as detecting unique solutions. \emph{Theoretical Computer Science, 47}, 85--93. Theorem~1.1, Corollary~1.2, p.~87, and Lemma~2.1, p.~88.
\bibitem[van der Geer and Katsura(2003)]{vdGK2003} van der Geer, G., \& Katsura, T. (2003). On the height of Calabi--Yau varieties in positive characteristic. \emph{Documenta Mathematica, 8}, 97--113. Corollary~2.3, p.~100. \url{https://doi.org/10.4171/DM/140}.
\bibitem[Vlasenko(2019)]{Vlasenko2019} Vlasenko, M. (2019). Formal groups and congruences. \emph{Transactions of the American Mathematical Society, 371}(2), 883--902. Theorem~2(i), arXiv:1509.06002v5, p.~3. \url{https://doi.org/10.1090/tran/7283}.
\bibitem[Wan(2008)]{Wan2008} Wan, D. (2008). Algorithmic theory of zeta functions over finite fields. In J. P. Buhler \& P. Stevenhagen (Eds.), \emph{Algorithmic number theory: Lattices, number fields, curves and cryptography} (pp.~551--578). Cambridge University Press. Theorem~6.3 and Corollary~6.4, p.~571. \url{https://library.slmath.org/books/Book44/files/17wan.pdf}.
\bibitem[Yobuko(2019)]{Yobuko2019} Yobuko, F. (2019). Quasi-Frobenius splitting and lifting of Calabi--Yau varieties in characteristic $p$. \emph{Mathematische Zeitschrift, 292}(1--2), 307--316. \url{https://doi.org/10.1007/s00209-018-2198-7}. Author manuscript: \url{https://arxiv.org/abs/1704.05604v2}.
\end{thebibliography}
\end{document}
